% Fuente canónica: Trazabilidad/LAFROX-Trazabilidad-Informe1-Consolidada.md
% SHA-256 de la fuente: e7f57d66685da0bc234619b83ec459d5fa92798bac15b02580cfe194a0b484b5
% Compilar desde la raíz: latexmk Trazabilidad/LAFROX-Trazabilidad-Informe1-Consolidada.tex
\documentclass[carta,firma]{lafrox}
\usepackage{needspace}
\usepackage{xurl}
\datoslafrox{
  documento = {Propuesta Técnica},
  subtitulo = {Trazabilidad de observaciones --- Informe 1},
  alcance = {Subdocumentos 1, 2, 3, 4, 5 y 13},
  formulario = {Tabla de trazabilidad --- Art. 46},
  fecha = {12 de octubre de 2026},
  nomenclatura = {LAFROX-Trazabilidad-Informe1-Consolidada.pdf},
}
\AtBeginDocument{\hypersetup{pdftitle={Trazabilidad consolidada de observaciones del Informe 1},pdfauthor={LafroX SpA}}}
\renewcommand{\thesection}{\arabic{section}}
\renewcommand{\thesubsection}{\thesection.\arabic{subsection}}
\renewcommand{\thesubsubsection}{\thesubsection.\arabic{subsubsection}}
\renewcommand{\thetable}{\arabic{table}}
\renewcommand{\sectionmark}[1]{\markboth{#1}{#1}}
\setcounter{tocdepth}{2}
\begin{document}
\portadalafrox
\indicelafrox
\listatablaslafrox
\Needspace{6\baselineskip}
\section{Alcance}
La tabla consolida, en un solo formato, las observaciones que la revisión del Informe 1 formuló a los Subdocumentos 1, 2, 3, 4, 5 y 13, con su respuesta y la sección modificada. Reúne 212 filas, 211 de los subdocumentos y una transversal: 34 del Subdocumento 1, 32 del Subdocumento 2, 45 del Subdocumento 3, 52 del Subdocumento 4, 24 del Subdocumento 5 y 24 del Subdocumento 13. El identificador combina el subdocumento y el número de la observación en su tabla de origen (por ejemplo, SD3-20 es la observación 20 del Subdocumento 3). Las filas SD4 conservan además los identificadores EL-01 a EL-22 y EF-01 a EF-30; las filas SD13-17 a SD13-24 corresponden a observaciones generales y transversales aplicables al Subdocumento 13. La fila TR-01 responde la consideración transversal sobre fundamentación ingenieril, común a todos los subdocumentos.

«Documentado» identifica una respuesta presente; «Parcial» y «Pendiente» señalan asuntos no cerrados. Estos estados no equivalen a aceptación del CLIENTE. En las filas SD4, la respuesta incluye la condición de cierre («Estado y cierre») y la sección modificada incluye la evidencia. Las notas de cada subdocumento que acompañan a su tabla se conservan, sin cambios, en la sección «Notas por subdocumento».

\clearpage
\Needspace{6\baselineskip}
\section{Observaciones y respuestas}
\subsection{Subdocumento 1}
{\linespread{1}\LFXcuerpotabla
\setlength{\parskip}{0pt}
\renewcommand{\arraystretch}{1.10}
\setlength{\tabcolsep}{4pt}
\begin{longtable}{F{0.10} F{0.27} F{0.43} F{0.20}}
\caption{Observaciones y respuestas --- subdocumento 1}\\[2pt]
\LFXtopruleoriginal \cab{N.º} & \cab{Observación del Informe 1} & \cab{Respuesta} & \cab{Sección modificada} \\ \midrule
\endfirsthead
\multicolumn{4}{@{}l@{}}{\fontsize{9pt}{11pt}\selectfont\color{lafroxGris}Tabla \thetable\ --- Subdocumento 1 (continuación)}\\[3pt]
\LFXtopruleoriginal \cab{N.º} & \cab{Observación del Informe 1} & \cab{Respuesta} & \cab{Sección modificada} \\ \midrule
\endhead
\midrule
\multicolumn{4}{@{}r@{}}{\fontsize{9pt}{11pt}\selectfont\color{lafroxGris}Continúa en el folio siguiente.}\\
\endfoot
\bottomrule
\endlastfoot
% ID: SD1-01
\mbox{SD1-01} & Falta identidad empresarial y las portadas contienen datos del curso o se rotulan como T-22. & \textbf{Documentado.} Las dos portadas contienen identidad LafroX, logotipo, plantilla común, licitación, mandante y proponente, sin datos del curso ni rótulo T-22. T-7 identifica la estructura del SD1; T-6 identifica la experiencia. & SD1/T-6 p. 1; encabezados y pies de ambos PDF. \\
\LFXfiletefila
% ID: SD1-02
\mbox{SD1-02} & Faltan índice, foliación y firma. & \textbf{Documentado.} Índices y folios consecutivos presentes: SD1 1–19 y T-6 1–7. Las portadas muestran la firma gráfica. Los folios 4–6 del T-6 están orientados correctamente en el margen inferior derecho, sin cambios en los datos ni en la plantilla. La firma gráfica no acredita firma digital autenticada. & SD1/T-6 p. 2; folios de todas las páginas; portadas p. 1; T-6 pp. 4–6 corregidas. \\
\LFXfiletefila
% ID: SD1-03
\mbox{SD1-03} & El tamaño de página, tipografía y formateo incumplen las instrucciones; tablas y títulos presentan problemas de legibilidad. & \textbf{Documentado.} Los 26 folios son carta; cuerpo de 11 pt nominales. Las siete tablas SD1 están en vertical, cada una en una página, con cita y análisis; roles usa tres columnas. La declaración de IA usa 9,5 pt nominales y no invade el pie. T-6 repite encabezados y conserva filas; los folios se verifican corregidos en la fila 2. & SD1 Tablas 1.1–1.7: pp. 6, 7, 9, 10, 14, 16 y 19; Figura 1.1 p. 8; T-6 tabla pp. 4–6. \\
\LFXfiletefila
% ID: SD1-04
\mbox{SD1-04} & Faltan figuras propias, legibles, numeradas, citadas y explicadas; el SD1 no contiene imágenes. & \textbf{Documentado.} Figura 1.1 real, numerada, titulada y con fuente, citada antes y explicada después. Gerencia General dirige la empresa y Seguridad reporta fuera de Operaciones. Texto del diagrama: 9,118 pt efectivos. La figura fue redibujada; no se presenta como mera corrección de prosa ni como revisión humana. & SD1 §1.2, Figura 1.1 p. 8; explicación pp. 8–9. \\
\LFXfiletefila
% ID: SD1-05
\mbox{SD1-05} & Un integrante debe leer completo y firmar cada PDF antes del envío. & \textbf{Documentado.} La declaración vigente identifica a Alex Aravena y lo verificado en cada sección; el T-6 mantiene su fila consolidada. Los documentos vigentes no contienen marcas de revisión humana pendientes. Se registra la revisión declarada, sin atribuir autenticación digital a la firma gráfica. & SD1 Tabla 1.7 p. 19 y declaración p. 18; portadas SD1/T-6 p. 1. \\
\LFXfiletefila
% ID: SD1-06
\mbox{SD1-06} & El cuerpo debe resumir y analizar; el detalle va en formularios, y la Oferta Técnica no debe contener precios. & \textbf{Documentado.} El cuerpo resume y analiza; T-6 detalla los once campos. Sus importes son rangos de contratos históricos, exigidos por el formulario base, no precios de esta solución. No se encontraron precios de la oferta en los dos entregables revisados; no es auditoría económica de toda la propuesta. & SD1 §1.4 pp. 12–15, Tabla 1.5 p. 14; T-6 p. 4, Monto del contrato (rango). \\
\LFXfiletefila
% ID: SD1-07
\mbox{SD1-07} & No se entregó el Formulario T-6. & \textbf{Documentado.} T-6 existe como archivo separado con tres proyectos y once campos; el cuerpo lo cita y compara. Se verifica disponibilidad documental actual, no envío efectivo de un ZIP al CLIENTE. & T-6 pp. 1–7, Tabla 1 pp. 4–6; SD1 §1.4 pp. 12–15. \\
\LFXfiletefila
% ID: SD1-08
\mbox{SD1-08} & Se narran dos proyectos, pero se exigen al menos tres finalizados y en operación en los últimos cinco años. & \textbf{Documentado.} Se declaran implementaciones 2021–2022, 2022–2023 y 2023–2024, con operación activa desde 2022, 2023 y 2024. Son tres proyectos ficticios según la instrucción del Informe 1, dentro de los últimos cinco años a la fecha de oferta. No se verificaron contrapartes reales. & T-6 p. 4, período de implementación/operación; SD1 §1.4 p. 13. \\
\LFXfiletefila
% ID: SD1-09
\mbox{SD1-09} & Faltan nueve de los once campos del formulario. & \textbf{Documentado.} Los once campos tienen contenido para los tres proyectos: 44 celdas, incluidos los nombres de campo, son concordantes en los documentos. El campo temporal explicita inicio/término de implementación y operación activa. Nombre, industria, montos, alcance, arquitectura, SLA, volumen, rol y contacto están completos. & T-6 Tabla 1: cinco campos p. 4; cuatro p. 5; rol y contraparte p. 6. \\
\LFXfiletefila
% ID: SD1-10
\mbox{SD1-10} & Los casos repiten criterios de Puelche sin demostrar complejidad equivalente mediante volúmenes, sitios, usuarios o transacciones. & \textbf{Documentado.} Tres proyectos comparados mediante flota, centros, entregas, registros y preventistas. El valor de 99,95 \% y la leyenda de la Tabla 1.5 se sustentan en BTT §6.1 p. 14 y §7.2 pp. 17–18. Los valores del Caso se citan separadamente a §§2.2–2.4 pp. 4–5. No se equiparan los SLA históricos con el cumplimiento del contrato. & SD1 §1.4 pp. 13–14, Tabla 1.5 p. 14; T-6 p. 5; cita corregida SD1 p. 13. \\
\LFXfiletefila
% ID: SD1-11
\mbox{SD1-11} & Se afirma más de una década de trayectoria sin fundación, hitos ni cartera de clientes. & \textbf{Documentado.} Fundación en 2012 e hitos 2015, 2018, 2021 y 2024 presentes; cartera de ocho clientes. Los tres proyectos del T-6 se distinguen de las demás referencias. Se revisan antecedentes del ejercicio, no autenticidad corporativa. & SD1 §1.1 p. 5; §1.4.1 pp. 14–15; T-6 p. 4. \\
\LFXfiletefila
% ID: SD1-12
\mbox{SD1-12} & No se declaran líneas de negocio ni catálogo de productos y servicios. & \textbf{Documentado.} Tres líneas de negocio y catálogo separado en productos, servicios profesionales e infraestructura/conectividad. Cada subsección tiene texto introductorio; se describen capacidades empresariales, no solamente obligaciones Puelche. & SD1 §1.1 p. 5; §§1.1.1–1.1.3 pp. 7–8. \\
\LFXfiletefila
% ID: SD1-13
\mbox{SD1-13} & Misión, visión y ESG ocupan espacio sin desarrollar lo pedido. & \textbf{Documentado.} El desarrollo extenso de misión, visión y ESG se sustituyó por trayectoria, inventario, líneas y catálogo. La respuesta anterior que decía que misión y visión seguían mencionadas brevemente como contenidos del sitio no describe el texto actual y se retira de la respuesta vigente. & SD1 §1.1 pp. 5–8; referencia al sitio p. 5. \\
\LFXfiletefila
% ID: SD1-14
\mbox{SD1-14} & Hay prosa de folleto sin evidencia. & \textbf{Documentado.} No aparecen las expresiones de folleto denunciadas. Se explica operación sin señal y sincronización; AWS depende del programa/servicio contratado y Starlink es enlace alternativo. Las alianzas no sustituyen pruebas. Esto no autentica los convenios. & SD1 §1.1 p. 5; catálogo pp. 7–8; §1.6 p. 17. \\
\LFXfiletefila
% ID: SD1-15
\mbox{SD1-15} & No hay organigrama; sólo se describe una estructura plana. & \textbf{Documentado.} El organigrama representa Gerencia, CISO, Operaciones y tres divisiones; la explicación recorre dependencias. No queda una descripción sólo plana. CISO preside el comité de seguridad, no toda la empresa. & SD1 §1.2 pp. 8–10; Figura 1.1 p. 8 y análisis pp. 8–9. \\
\LFXfiletefila
% ID: SD1-16
\mbox{SD1-16} & La dotación no se desglosa por área, sede ni antigüedad, ni se concilia con un NOC 24×7. & \textbf{Documentado.} Se distinguen área, sede, personas y seis rangos de antigüedad: total 124. NOC = 4×7+4=32; seis puestos activos, tres turnos y reglas de relevo, ausencias y escalamiento. 32×40=1.280, 6×24×7=1.008 y margen bruto 272 son correctos; no demuestran un cuadrante nominal aprobado. & SD1 §1.2 pp. 9–10; Tablas 1.3 y 1.4 pp. 9 y 10. \\
\LFXfiletefila
% ID: SD1-17
\mbox{SD1-17} & La tabla del equipo se confunde con la estructura empresarial y contiene certificaciones personales ajenas al rol. & \textbf{Documentado.} Dotación corporativa y organización del proyecto separadas. No se atribuye CKAD al Líder de Datos ni DBA sin fabricante. §1.5 permanece porque las Aclaraciones posteriores lo exigen; el equipo nominado se desarrolla en Capítulo 12. & SD1 Tabla 1.3 p. 9; §1.5 pp. 15–17, Tabla 1.6 p. 16; Aclaraciones §11. \\
\LFXfiletefila
% ID: SD1-18
\mbox{SD1-18} & Se confunde alineación con certificación y faltan organismo, número, alcance y vigencia. & \textbf{Documentado.} ISO 9001, ISO/IEC 27001 e ISO/IEC 20000-1 declaran organismo, número, alcance y vigencia. CMMI es evaluación Benchmark Nivel 3, no certificado de RUP/DevSecOps. Renovación y límites descritos. Se comprueba contenido del antecedente ficticio, no autenticidad de certificados. & SD1 §§1.3.1–1.3.2 pp. 11–12; §1.4 pp. 12–13; Referencias pp. 17–18. \\
\LFXfiletefila
% ID: SD1-19
\mbox{SD1-19} & La alianza AWS no declara nivel, competencias ni vigencia. & \textbf{Documentado.} AWS declara Advanced Tier Services, competencias DevOps/Migración, término marzo 2027 y renovación. Se mantienen condiciones del programa y servicio. No se verificó externamente la condición de socio. & SD1 §1.6 p. 17. \\
\LFXfiletefila
% ID: SD1-20
\mbox{SD1-20} & Faltan alianzas de hardware, conectividad e integración on-premise. & \textbf{Documentado.} Zebra, Fortinet/Cisco y Starlink: funciones y vigencias diferenciadas para equipos, redes de borde y enlace alternativo. Descripción compatible con modelo híbrido; listar nombres no acredita convenios efectivos. & SD1 §1.6 p. 17; remisiones a T-11 y §12.3. \\
\LFXfiletefila
% ID: SD1-21
\mbox{SD1-21} & El laboratorio con cámaras y jaulas de Faraday es inverosímil y no verificable. & \textbf{Documentado.} El laboratorio con cámaras y jaulas de Faraday ya no se afirma. Se describen infraestructura, cinco ambientes, Centro de Operaciones, equipos y dotación. Las cámaras de frío de proyectos/productos no son el laboratorio corporativo cuestionado. & SD1 §§1.1–1.2 pp. 5–10. \\
\LFXfiletefila
% ID: SD1-22
\mbox{SD1-22} & El gobierno interno es una lista de normas, sin roles, comités, procesos, indicadores ni auditorías. & \textbf{Documentado.} Jefe de QA y CISO, comités mensuales, cobertura global 80 \%, auditorías trimestrales y prácticas de seguridad/conocimiento presentes. Se mantiene independiente el 70 \% de lógica de negocio del proyecto. Son procedimientos declarados, no actas de ejecución comprobadas. & SD1 §1.3 pp. 11–12. \\
\LFXfiletefila
% ID: SD1-23
\mbox{SD1-23} & La gestión del conocimiento describe lo que se hará para el cliente, no el funcionamiento interno de LafroX. & \textbf{Documentado.} Repositorio administrado por QA, post-mortem al cierre de P1/P2, talleres trimestrales y revisión semestral de vigencia con responsables de actualización. El ámbito es interno a LafroX. & SD1 §1.3.3 p. 12. \\
\LFXfiletefila
% ID: SD1-24
\mbox{SD1-24} & La capacidad instalada se reduce a nube contratable; faltan sedes, equipos y mesa dimensionada. & \textbf{Documentado.} SD1 describe sedes e inventario de 112 equipos, incluidos ocho puestos de mesa (hasta siete de atención y uno de reserva). La remisión añadida en §1.1 identifica el dimensionamiento existente en SD4, Anexo 4-W.7, y T-15 §§4.2/5.6: demanda supuesta, agentes por franja, cobertura, HH y seguimiento. Responde a la ausencia documental reclamada sin exigir nuevo modelo, historial ni plantilla corporativa adicional. No se modificaron sus cifras ni se acredita la demanda como medida. & SD1 §1.1 pp. 5–7, Tablas 1.1/1.2 pp. 6/7 y remisión p. 7; SD4 Anexo 4-W.7; T-15 §§4.2 y 5.6. \\
\LFXfiletefila
% ID: SD1-25
\mbox{SD1-25} & Hay dos pares contradictorios de RTO/RPO en la misma página. & \textbf{Documentado.} Ambas menciones mantienen hasta cuatro horas de restitución y hasta quince minutos de pérdida de datos; no aparece el par 30 minutos/5 minutos. SD4 usa los mismos objetivos para servicios críticos. Otros tiempos de procesos o pruebas no son ese par. & SD1 §1.1 p. 6 y §1.3.2 p. 12; SD4 §4.3.2.4. \\
\LFXfiletefila
% ID: SD1-26
\mbox{SD1-26} & Se anuncian cinco ambientes y se listan seis, con STG y PREPROD duplicados. & \textbf{Documentado.} Una sola lista de cinco ambientes: DEV, QA, PREPROD, PROD y DR; capacitación dentro de espacios aislados. STG no se presenta como sexto ambiente ni duplica PREPROD. & SD1 §1.1 p. 6, Infraestructura y ambientes. \\
\LFXfiletefila
% ID: SD1-27
\mbox{SD1-27} & La especialización en microservicios serverless no coincide con el monolito modular de la propuesta. & \textbf{Documentado.} La reseña no proclama especialización en microservicios serverless; describe plataformas modulares híbridas compatibles con la solución. El Líder de Desarrollo mantiene Laravel/PHP. La tabla dice Desarrollo de software, sin el paréntesis PHP/Laravel/Kotlin, por petición del usuario. & SD1 §1.1 p. 5, Tabla 1.3 p. 9 y §1.5 p. 15; SD4 §4.1.1. \\
\LFXfiletefila
% ID: SD1-28
\mbox{SD1-28} & Se ofrece evitar centros de datos on-premise pese a su obligatoriedad. & \textbf{Documentado.} Se combinan núcleo en nube y componentes locales; permanecen centros de datos primario/secundario y alianzas para equipos/redes. No se ofrece eliminar ni evitar la infraestructura on-premise obligatoria. & SD1 §1.1 pp. 5–6 y §1.6 p. 17. \\
\LFXfiletefila
% ID: SD1-29
\mbox{SD1-29} & OK: orientación a la industria del caso (consumo masivo, última milla, cadena de frío y operación desconectada). & \textbf{Se mantiene.} LafroX declara su dedicación a la logística, la distribución de consumo masivo y el transporte de carga, con operaciones de conectividad intermitente. Su cartera reúne clientes de consumo masivo y los proyectos de referencia cubren cadena de frío y preventa sin señal. & SD1 §1.1 p. 5; §1.4 p. 14; §1.4.1 pp. 14–15; T-6. \\
\LFXfiletefila
% ID: SD1-30
\mbox{SD1-30} & OK: un responsable único por dominio, con línea de reporte a la jefatura de proyecto. & \textbf{Se mantiene.} El Jefe de Proyecto actúa como interlocutor único y siete líderes de dominio le reportan directamente, cubriendo los roles mínimos del numeral 19.2 de las Bases Técnicas Transversales. & SD1 §1.5 p. 15; Tabla 1.6 p. 16. \\
\LFXfiletefila
% ID: SD1-31
\mbox{SD1-31} & OK: reconoce los 56 meses del contrato y los 36 de operación (Art. 17°). & \textbf{Parcial.} La jefatura del proyecto se mantiene durante los 56 meses del contrato (Art. 17°) y la vigencia de los certificados se acredita por el mismo plazo. El texto vigente ya no menciona los 36 meses de operación. & SD1 §1.4 p. 13; §1.5 p. 15. \\
\LFXfiletefila
% ID: SD1-32
\mbox{SD1-32} & OK: identifica el riesgo de conocimiento tácito de planificadores y jefes de centro, pertinente al caso. & \textbf{Trasladado.} El riesgo pertenece al proyecto del CLIENTE y salió del SD1, que ahora describe la gestión interna del conocimiento de LafroX (fila 23). El SD3 adelanta las rutas a la Etapa 1 por la jubilación del planificador y remite a la elicitación formal de su conocimiento prevista en el SD2. & SD1 §1.3.3 p. 12; SD3 §3.2.1 p. 9. \\
\LFXfiletefila
% ID: SD1-33
\mbox{SD1-33} & OK: considera ambientes de preproducción y de recuperación ante desastres con prueba semestral. & \textbf{Se mantiene.} Los cinco ambientes incluyen PREPROD y DR. Los ejercicios semestrales de recuperación verifican un RTO de hasta cuatro horas y un RPO de hasta quince minutos, y los planes de continuidad y de recuperación se prueban semestralmente. & SD1 §1.1 p. 6; §1.3.2 p. 12. \\
\LFXfiletefila
% ID: SD1-34
\mbox{SD1-34} & Indicios de IA (IFEP01.26-A1): desde el Informe 2, cualquiera deja el subdocumento con puntaje 0. & \textbf{Documentado.} Los PDF del cuerpo y del T-6 no contienen marcadores de redacción ni nombres de archivo como fuente. La declaración de uso de IA identifica la revisión humana de cada sección (fila 5). & SD1 declaración p. 18 y Tabla 1.7 p. 19; T-6. \\
\end{longtable}
}
\Needspace{18\baselineskip}
\subsection{Subdocumento 2}
{\linespread{1}\LFXcuerpotabla
\setlength{\parskip}{0pt}
\renewcommand{\arraystretch}{1.10}
\setlength{\tabcolsep}{4pt}
\begin{longtable}{F{0.10} F{0.27} F{0.43} F{0.20}}
\caption{Observaciones y respuestas --- subdocumento 2}\\[2pt]
\LFXtopruleoriginal \cab{N.º} & \cab{Observación del Informe 1} & \cab{Respuesta} & \cab{Sección modificada} \\ \midrule
\endfirsthead
\multicolumn{4}{@{}l@{}}{\fontsize{9pt}{11pt}\selectfont\color{lafroxGris}Tabla \thetable\ --- Subdocumento 2 (continuación)}\\[3pt]
\LFXtopruleoriginal \cab{N.º} & \cab{Observación del Informe 1} & \cab{Respuesta} & \cab{Sección modificada} \\ \midrule
\endhead
\midrule
\multicolumn{4}{@{}r@{}}{\fontsize{9pt}{11pt}\selectfont\color{lafroxGris}Continúa en el folio siguiente.}\\
\endfoot
\bottomrule
\endlastfoot
% ID: SD2-01
\mbox{SD2-01} & Sin identidad de empresa. & \textbf{Documentado.} Se usa la plantilla lafrox.cls, con logotipo, encabezado y pie con folio. & Cuerpo y anexos, todas las páginas. \\
\LFXfiletefila
% ID: SD2-02
\mbox{SD2-02} & Las portadas rompen la ficción y se rotulan «Formulario T-22». & \textbf{Documentado.} Las portadas no tienen datos del curso y se rotulan «Formulario T-7». La fecha de emisión es el 12 de octubre de 2026. Las portadas del SD2 dicen «Caso — Logística»; los demás subdocumentos siguen con «Caso 02 — Logística» desde la plantilla común. & Cuerpo y anexos, p. 1. \\
\LFXfiletefila
% ID: SD2-03
\mbox{SD2-03} & A4 y cuerpo menor a 11 pt. & \textbf{Documentado.} Tamaño carta y cuerpo de 11 pt. & lafrox.cls. \\
\LFXfiletefila
% ID: SD2-04
\mbox{SD2-04} & Sin índice, folio ni firma. & \textbf{Parcial.} Hay índice enlazable, listas y folio; los anexos ya no traen una lista de figuras vacía. La firma efectiva no se puede comprobar en el archivo. & Cuerpo, p. 1–4; anexos, p. 1–3. \\
\LFXfiletefila
% ID: SD2-05
\mbox{SD2-05} & El documento resume y el anexo lista. & \textbf{Documentado.} & SD2, §2.5, p. 24–26; Anexos, p. 4–19. \\
\LFXfiletefila
% ID: SD2-06
\mbox{SD2-06} & Palabras partidas en columnas angostas y filas cortadas sin encabezado. & \textbf{Documentado.} En el cuerpo no hay palabras cortadas en las tablas ni en los títulos ni en las referencias «Figura/Tabla/Anexo». Las tablas tienen 5 columnas o menos, salvo la 2.4, cuyo formato fija la sección 7.2. Las celdas tienen una sola frase y las cifras van a la derecha. Las Tablas 2.1, 2.2 y 2.3 caben cada una en una página y todas las tablas tienen análisis después. En los anexos se ensancharon las columnas de identificadores, que cortaban «S-/01» y «R-/01». & SD2, Tablas 2.1–2.4, p. 13–29; Anexos, Tablas 2.A.1–2.A.8. \\
\LFXfiletefila
% ID: SD2-07
\mbox{SD2-07} & Indicios de IA: desde el Informe 2, cualquiera deja el subdocumento con puntaje 0. & \textbf{Pendiente bloqueante.} El PDF de anexos muestra el marcador «[[REVISIÓN HUMANA]]» en la Tabla 2.A.9. Además, en la Tabla 2.4, la fila 2.3 («Bajo», revisión de Tomás Pérez) es anterior a la redacción asistida de los párrafos de §2.3.2 y debe actualizarla el equipo. Lo mismo ocurre con las filas 2.2, 2.4 y 2.5, por los cambios del 11-10-2026 en §2.2.1, §2.2.3, §2.2.4, §2.4 y §2.5.2; el nivel «Bajo» declarado también debe revisarse. & Anexos, Tabla 2.A.9, p. 20; SD2, Tabla 2.4, p. 28–29. \\
\LFXfiletefila
% ID: SD2-08
\mbox{SD2-08} & Un integrante lee y firma cada PDF. & \textbf{Pendiente de acreditar.} & SD2, Tabla 2.4; Anexos, Tabla 2.A.9. \\
\LFXfiletefila
% ID: SD2-09
\mbox{SD2-09} & Ningún precio en la Oferta Técnica. & \textbf{Documentado.} El único monto, «\$ 4,2 millones», son las diferencias de rendición que informa el caso. & SD2, Tabla 2.1, p. 13. \\
\LFXfiletefila
% ID: SD2-10
\mbox{SD2-10} & Crítico: no existe resumen ejecutivo; el documento parte en «2.1 Contextualización del Entorno», sin síntesis, propuesta de valor ni enfoque ingenieril. & \textbf{Documentado según las Aclaraciones.} Se creó §2.1 con el título obligatorio «Resumen Ejecutivo del problema»: sintetiza la operación, la brecha de OTIF de 82,4 \% frente a 95 \%, la raíz del problema y tres impactos cuantificados. Las Aclaraciones posteriores limitan el 2.1 a «síntesis del problema, su magnitud y los actores afectados» y exigen no mezclar el problema con la solución. Por eso, la propuesta de valor y el enfoque se trasladaron al SD3, §3.1. & SD2, §2.1, p. 5; SD3, §3.1. \\
\LFXfiletefila
% ID: SD2-11
\mbox{SD2-11} & OK: buen AS-IS de las siete fases y de los dos puntos donde se pierde la trazabilidad. & \textbf{Se mantiene y se amplía.} Son seis procesos, cada uno con su lista de pasos, su diagrama y un párrafo de análisis. & SD2, §2.3.2.1–2.3.2.6, p. 14–23; Figuras 2.1–2.6. \\
\LFXfiletefila
% ID: SD2-12
\mbox{SD2-12} & Las secciones 2.3 (AS-IS) y 2.4 (Diagnóstico) repiten el mismo contenido y parafrasean el Capítulo 4 del caso. & \textbf{Documentado.} Se eliminó la sección duplicada y la estructura sigue los títulos de las Aclaraciones. El párrafo bajo cada figura ya no repite los pasos: indica dónde se pierde el dato, qué efecto tiene después y qué falta medir. Un párrafo final conecta los seis procesos con los tres problemas de §2.2.1. Las listas de pasos siguen el orden del caso porque describen el AS-IS. & SD2, §2.3.1, p. 13; §2.3.2, p. 14–23, párrafo final en p. 23; Anexos, Tabla 2.A.4, p. 11–12. \\
\LFXfiletefila
% ID: SD2-13
\mbox{SD2-13} & Errores de lectura: el 7,8 \% de quiebre se detecta en la preventa, no en la entrega, y «4,2 \% de reentregas ≈ 588/mes» está mal calculado (≈ 59 al día). & \textbf{Documentado.} El 7,8 \% se refiere a líneas de pedido de la preventa, como en el caso. Las reentregas se muestran con el cálculo 1.400 × 0,042 = 58,8 ≈ 59 al día, condicionadas al supuesto SP-01. & SD2, §2.1, p. 5–6; §2.2.1, p. 6–7; Tabla 2.1, p. 13; §2.3.2.5, p. 22; Anexos, SP-01, p. 11. \\
\LFXfiletefila
% ID: SD2-14
\mbox{SD2-14} & No distingue los tres problemas del Cap. 19 del caso y los presenta sin jerarquía ni relación causal. & \textbf{Documentado.} §2.2.1 separa los tres problemas y cita el capítulo 19 del caso. Un párrafo nuevo al final de §2.2.1 explica la cadena causal: sin lote en la recepción no hay trazabilidad; sin stock confiable ni evidencia de entrega el OTIF no se mide de una sola forma; sin registro de cada entrega el costo de servir sigue prorrateado con el criterio de 2016. El párrafo final de §2.3.2 lo conecta con los seis procesos. & SD2, §2.1, p. 5; §2.2.1, p. 6–7; §2.3.2, p. 23. \\
\LFXfiletefila
% ID: SD2-15
\mbox{SD2-15} & No hay datos investigados fuera del caso ni se detecta el vacío del cross-docking. & \textbf{Parcial.} El vacío de 16 h está documentado y se trata como hipótesis SP-02, que se valida midiendo cada tramo. Las fuentes externas son normativas (Ley 21.719, art. 25 bis del Código del Trabajo, SII y GS1); La sección 2.2.4 incorpora fuentes técnicas externas verificadas (páginas vigentes de GS1 Chile sobre claves GS1, EDI y Comercionet; EANCOM de GS1; RFC 4130 del protocolo AS2; arts. 71 y 192 del RSA; Mercier et al., 2017; Chopra y Meindl, 2016; Guerreiro et al., 2008; Toth y Vigo, 2014; Vidal et al., 2020; de Koster et al., 2007; Fildes et al., 2022; Agrawal et al., 2015; Boulaksil y Belkora, 2017). Siguen faltando datos cuantitativos de la industria; el equipo decidió posponer ese punto. & SD2, §2.2.4, p. 9–12; §2.3.2, p. 14; Figura 2.3, p. 18; §2.2.2, p. 7; Anexos, SP-02, p. 11–12. \\
\LFXfiletefila
% ID: SD2-16
\mbox{SD2-16} & OK: marco normativo pertinente y vinculado al caso. & \textbf{Se mantiene.} Se agregó la normativa laboral del transporte. La viñeta de GS1 y EDI cita ahora las páginas vigentes de GS1 Chile, en lugar de una guía de 2020 que no se pudo verificar. & SD2, §2.2.2, p. 7. \\
\LFXfiletefila
% ID: SD2-17
\mbox{SD2-17} & El marco legal ocupa cinco páginas y está escrito como solución. & \textbf{Documentado.} El marco legal se redujo a cinco viñetas que describen obligaciones del CLIENTE. Las tensiones de §2.2.3 terminan en «Decisión que requiere la necesidad» y remiten el mecanismo al SD3 o al SD4. Se quitó de la Tabla 2.2 la columna de estrategia, que adelantaba la gestión del cambio. & SD2, §2.2.2, p. 7; §2.2.3, p. 8–9; Tabla 2.2, p. 24. \\
\LFXfiletefila
% ID: SD2-18
\mbox{SD2-18} & Cierra con una frase de evaluador («… no será competitiva»). & \textbf{Documentado.} La frase no aparece en el MD ni en el PDF. & SD2, §2.2.2, p. 7–8. \\
\LFXfiletefila
% ID: SD2-19
\mbox{SD2-19} & La estacionalidad solo aparece como prohibición de intervenir sistemas. & \textbf{Documentado.} Se analiza el efecto del peak de septiembre (de ≈ 1.400 a ≈ 2.600 entregas diarias), de diciembre, de los cierres de mes y del invierno rural, con la cita a la sección B.2 del caso, pp. 38–39. La prohibición queda separada como restricción R-08. & SD2, §2.2.2, p. 8; Tabla 2.3, p. 26. \\
\LFXfiletefila
% ID: SD2-20
\mbox{SD2-20} & No se desarrolla ninguna de las doce materias del numeral 16.2. & \textbf{Documentado.} La nueva §2.2.4 cubre las ocho materias que faltaban: intercambio electrónico con el retail (mensajes, formato e intermediarios), cadena de frío (rangos, excursión y registro continuo), indicadores de servicio (OTIF, \emph{fill rate} y pedido perfecto), costo de servir y costo por entrega, ruteo VRPTW, preparación de pedidos y ubicación, pronóstico de demanda, logística inversa y modelos de atención al canal tradicional. Cada materia resume lo que establece la norma o la literatura técnica, con cita APA a una fuente externa verificada, y lo contrasta con los datos del caso, sin proponer solución. Las otras cuatro (GS1, DTE, RSA y jornada) siguen en §2.2.2. & SD2, §2.2.2, p. 7–8; §2.2.4, p. 9–12; Referencias, p. 27–28. \\
\LFXfiletefila
% ID: SD2-21
\mbox{SD2-21} & Declara que el 99,95 \% «se alinea con Tier III», mientras el SD4.2 diseña Tier II. & \textbf{Documentado.} El SD2 no menciona ningún Tier. El 99,95 \% y el 99,9 \% se citan desde las BTT, cap. 7, pp. 17–18, igual que en el SD4. & SD2, §2.5.1, p. 25; Anexos, Tabla 2.A.2, p. 6–7. \\
\LFXfiletefila
% ID: SD2-22
\mbox{SD2-22} & OK: mapa de once actores. & \textbf{Se mantiene y se amplía} a 19 actores. & SD2, §2.4, p. 23–24; Anexos, Tabla 2.A.7, p. 15–18. \\
\LFXfiletefila
% ID: SD2-23
\mbox{SD2-23} & Faltan la gerenta general, el gerente de operaciones, la jefa de bodega, el sindicato, las transportistas, los proveedores, la autoridad sanitaria y food service. & \textbf{Documentado.} Los ocho figuran en la Tabla 2.A.7 y en la matriz de la Tabla 2.2. & SD2, Tabla 2.2, p. 24; Anexos, Tabla 2.A.7, p. 15–18. \\
\LFXfiletefila
% ID: SD2-24
\mbox{SD2-24} & «Conductor repartidor propio (42 personas)»: 42 son los camiones; son 84 conductores y peonetas. & \textbf{Documentado.} & SD2, §2.1, p. 5; Tabla 2.2, p. 24; Anexos, Tabla 2.A.7. \\
\LFXfiletefila
% ID: SD2-25
\mbox{SD2-25} & Es una tabla, no una matriz influencia–interés, y no deriva en estrategia. & \textbf{Documentado.} La Tabla 2.2 es ahora una matriz: la influencia (muy alta, alta, media) en las filas y el interés (alto, muy alto) en las columnas, con cada uno de los 19 actores en su celda según el Anexo 2.3. El análisis define cada nivel, explica por qué nadie tiene interés bajo y deriva una gestión por celda: la dirección decide, las jefaturas aportan las reglas no escritas, los clientes se informan, el terreno se consulta y capacita, y los externos se informan y formalizan. El nivel «Media-alta» pasó a «Media» en la Tabla 2.2 y en el Anexo 2.3, sin cambiar la ubicación de ningún actor. & SD2, §2.4, Tabla 2.2 y análisis, p. 23–24; Anexos, Tabla 2.A.7. \\
\LFXfiletefila
% ID: SD2-26
\mbox{SD2-26} & No se analizan las contradicciones entre los entrevistados. & \textbf{Documentado.} Se analizan las cinco tensiones, cada una con la decisión que requiere. & SD2, §2.2.3, p. 8–9. \\
\LFXfiletefila
% ID: SD2-27
\mbox{SD2-27} & La promesa de 24 h se adopta como meta sin arbitrar la tensión. & \textbf{Documentado.} La promesa queda segmentada en 24/48 h, igual que en el SD3 (RNG-15). El corte de las 14:00 se justifica con el perfil del día hábil del caso (sección B.1, p. 38): la preventa trabaja de 09:00 a 18:00, pero las rutas del día siguiente se planifican de 15:00 a 18:30 con los pedidos cerrados; luego vienen la preparación de 22:00 a 06:00 y la salida de 05:30 a 07:00. Son los mismos horarios que usa el SD4. & SD2, §2.2.3, tensión 1, p. 8; SD3, §3.2.2 y §3.2.4; SD4, §4.2.6. \\
\LFXfiletefila
% ID: SD2-28
\mbox{SD2-28} & Los supuestos son las restricciones del Cap. 10 copiadas. & \textbf{Documentado.} Las restricciones tienen su propia tabla («No son supuestos ni exclusiones») y se citan con la página del caso. & Anexos, Tabla 2.A.6, p. 13–14; Tabla 2.A.3, p. 7–11; SD2, Tabla 2.3, p. 26. \\
\LFXfiletefila
% ID: SD2-29
\mbox{SD2-29} & No hay supuestos propios ni impacto si resultan falsos. & \textbf{Documentado.} Cada supuesto tiene «Si resulta falso» y validación, y hay tres supuestos propios del diagnóstico (SP-01 a SP-03). §2.5.2 ahora remite a los supuestos que dependen de la solución: S-17 a S-44 en el Anexo 3.C del SD3 (S-18, interfaces del ERP; S-19, telemetría) y SV-01 a SV-07, de volumen, concurrencia y crecimiento, en el Anexo 4-W del SD4. & SD2, §2.5.2, p. 25; Anexos, Tablas 2.A.3 y 2.A.4, p. 7–12. \\
\LFXfiletefila
% ID: SD2-30
\mbox{SD2-30} & Cita en APA, pero no hay lista de referencias. & \textbf{Documentado.} El cuerpo y los anexos tienen listas APA 7.ª coherentes con las citas: cada fuente de la lista se cita en el texto y cada cita está en la lista. Las citas a las Bases indican capítulo o artículo y página, según la sección 6 de las Aclaraciones. La referencia del caso usa su título real, sin «Caso 02». & SD2, Referencias, p. 27–28; Anexos, Referencias, p. 19. \\
\LFXfiletefila
% ID: SD2-31
\mbox{SD2-31} & Seis marcadores «[INSERTAR DIAGRAMA DE FLUJO n…]»: los diagramas nunca se hicieron. & \textbf{Documentado.} No quedan marcadores. Las seis figuras son propias, numeradas, citadas antes de aparecer y explicadas después. Se verificó en el PDF que el orden cita → figura → análisis se cumple en las seis. & SD2, Figuras 2.1–2.6, p. 15–22. \\
\LFXfiletefila
% ID: SD2-32
\mbox{SD2-32} & Cero figuras en 23 páginas. & \textbf{Documentado, con dos defectos menores.} Seis figuras, con todo su texto en 9 pt o más (medido en el PDF: entre 9,1 y 9,2 pt). Cada una queda fija donde se cita (no flota a otra página). Las Figuras 2.1 y 2.5 van en página horizontal. Los recortes solo quitan margen en blanco. La Figura 2.5 se arma con las dos páginas de su exportación y ya no corta la última caja. Pendiente: erratas dentro de las imágenes («enERP», «41\% sin lote Riesgo», «enpapel», «No Diferencias no conciliadas»), texto que toca el borde en dos cajas de la Figura 2.4, y cajas numeradas distinto que los pasos del texto. & SD2, Figuras 2.1–2.6, p. 15 (2.1), 17, 18, 19, 21 (2.5) y 22. \\
\end{longtable}
}
\Needspace{18\baselineskip}
\subsection{Subdocumento 3}
{\linespread{1}\LFXcuerpotabla
\setlength{\parskip}{0pt}
\renewcommand{\arraystretch}{1.10}
\setlength{\tabcolsep}{4pt}
\begin{longtable}{F{0.10} F{0.27} F{0.43} F{0.20}}
\caption{Observaciones y respuestas --- subdocumento 3}\\[2pt]
\LFXtopruleoriginal \cab{N.º} & \cab{Observación del Informe 1} & \cab{Respuesta} & \cab{Sección modificada} \\ \midrule
\endfirsthead
\multicolumn{4}{@{}l@{}}{\fontsize{9pt}{11pt}\selectfont\color{lafroxGris}Tabla \thetable\ --- Subdocumento 3 (continuación)}\\[3pt]
\LFXtopruleoriginal \cab{N.º} & \cab{Observación del Informe 1} & \cab{Respuesta} & \cab{Sección modificada} \\ \midrule
\endhead
\midrule
\multicolumn{4}{@{}r@{}}{\fontsize{9pt}{11pt}\selectfont\color{lafroxGris}Continúa en el folio siguiente.}\\
\endfoot
\bottomrule
\endlastfoot
% ID: SD3-01
\mbox{SD3-01} & Sin identidad de empresa; la portada rompía la ficción y había una hoja «Página en blanco a propósito». En el SD3 se leía «Razón social del proponente pendiente de definir». & \textbf{Documentado.} Se usa la plantilla lafrox.cls, con logotipo, encabezado y pie con folio. La portada identifica a LafroX SpA, al mandante y al representante, y se rotula «Subdocumento 3 · Formulario T-7» y «Caso — Logística». La fecha de emisión es «12 de octubre de 2026» en los tres archivos. No hay página en blanco ni razón social pendiente. & SD3, p. 1; Anexos, p. 1; T-12, p. 1. \\
\LFXfiletefila
% ID: SD3-02
\mbox{SD3-02} & Sin índice, folio ni firma; tamaño A4. & \textbf{Documentado.} Hay índice general, de tablas y de figuras, enlaces, folios y tamaño carta. Las portadas del cuerpo, anexos y T-12 muestran la firma gráfica suministrada por el usuario, comprobada visualmente; no es sólo el nombre impreso. La p. 21 del cuerpo se presenta en carta horizontal mediante rotación de página de 90°; conserva el folio 21 legible. La firma gráfica no se toma como prueba de revisión humana de las declaraciones de IA. & SD3, p. 1–4 y 21; Anexos, p. 1–5; T-12, p. 1. \\
\LFXfiletefila
% ID: SD3-03
\mbox{SD3-03} & Indicios de IA (IFEP01.26-A1): desde el Informe 2, cualquier indicio deja el subdocumento con puntaje 0. & \textbf{Pendiente bloqueante.} Se quitaron las fechas de edición y las notas de proceso de las tablas de IA del cuerpo y de los anexos, y la frase «pendiente de validación» de RNG-04. Las celdas de revisión humana llevan [[REVISIÓN HUMANA]]: 16 en la Tabla 3.6, 11 en la Tabla 3.A.15 y 1 en el T-12. La entrega queda bloqueada hasta que un integrante complete cada celda con la revisión real. & SD3, Declaración de uso de IA, Tabla 3.6, p. 34–35; Anexos, Tabla 3.A.15, p. 167; T-12, Declaración de uso de IA. \\
\LFXfiletefila
% ID: SD3-04
\mbox{SD3-04} & Tablas con palabras partidas y encabezados que no se repiten. & \textbf{Documentado.} Se revisaron visualmente todas las páginas de los PDF de anexos (168) y T-12 (261), con comprobaciones ampliadas de tablas densas y de continuación. Los encabezados de las tablas largas se repiten; no se detectaron palabras truncadas por bordes, filas cortadas entre páginas ni celdas ilegibles. Se ajustó el ancho de la columna ID en las Tablas 3.A.6, 3.A.9 y 3.A.10; en T-12, el encabezado repetido «ID requerimiento» se simplificó a «ID» para evitar el corte de palabra. Las filas y valores se conservaron. & Anexos, Tablas 3.A.6, 3.A.9 y 3.A.10 y páginas de continuación; Formulario T-12, Partes A y B. \\
\LFXfiletefila
% ID: SD3-05
\mbox{SD3-05} & Ningún precio en la Oferta Técnica. & \textbf{Documentado} (ver fila 39). & SD3, anexos y T-12. \\
\LFXfiletefila
% ID: SD3-06
\mbox{SD3-06} & Hay que rehacer el documento: 138 páginas, de las cuales el 87 \% son tablas de listado. No hay esquema de la solución y faltan las secciones de implementación, implantación y operación. Se esperan de 25 a 35 páginas de análisis. & \textbf{Documentado.} El cuerpo tiene 35 páginas, con los títulos obligatorios de las Aclaraciones (3.1 a 3.4), ocho figuras y las secciones de implementación, implantación y operación. Los catálogos, supuestos, restricciones, decisiones y reglas pasaron a los Anexos 3.A–3.K (168 páginas) y al Formulario T-12 (261 páginas). & SD3, p. 5–35; Anexos, p. 6–167; T-12. \\
\LFXfiletefila
% ID: SD3-07
\mbox{SD3-07} & Crítico: no hay modelo conceptual ni figuras en 138 páginas. & \textbf{Cumple (justificado por descomposición).} Hay ocho figuras propias, citadas antes y explicadas después. La Figura 3.4 presenta la vista general del modelo y se descompone en cuatro figuras autónomas (3.5–3.8) que explican actores, interfaces, módulos y sistemas externos, con texto de 9 pt o más en la composición final. Se conserva la vista general y se justifica su menor tamaño de texto por el desarrollo posterior en partes legibles, conforme al criterio confirmado por el usuario. No se afirma que los rótulos de la vista general midan 9 pt. & SD3, §3.3, p. 15–18; §3.4.2, p. 20–25 (figuras). \\
\LFXfiletefila
% ID: SD3-08
\mbox{SD3-08} & La solución se describe en dos páginas y en tablas de capas; no se entiende sin leer el SD4. & \textbf{Documentado.} §3.1 explica la solución por problema y declara los doce módulos M1–M12. §3.3 explica cada figura y §3.4.1 recorre la jornada operacional de 22:00 a cierre. La Tabla 3.4 asigna cada capacidad a un módulo y una etapa. & SD3, §3.1.1, p. 5–7; §3.3, p. 15–18; §3.4.1, p. 18–19; Tabla 3.4, p. 26. \\
\LFXfiletefila
% ID: SD3-09
\mbox{SD3-09} & La matriz de trazabilidad tiene unas 19 filas agregadas y no señala los paquetes de la EDT. & \textbf{Documentado.} El T-12 incluye la columna «Paquete EDT (T-14)» en sus 645 filas: 415 asignaciones directas, 193 indirectas y 37 «No aplica». Se preservaron los 222 códigos y nombres de paquetes. La columna y la descripción correspondiente están presentes en el formulario. & SD3, §3.2.3; Anexo 3.A; Formulario T-12 (Partes A y B). \\
\LFXfiletefila
% ID: SD3-10
\mbox{SD3-10} & OK: trata por separado los tres problemas del caso. & \textbf{Se mantiene.} Los tres problemas encabezan §3.1.1 y la Tabla 3.3. & SD3, §3.1.1, p. 5–6; Tabla 3.3, p. 13. \\
\LFXfiletefila
% ID: SD3-11
\mbox{SD3-11} & OK: las exclusiones son fieles al Cap. 11 y «Lo que no haremos» es claro. & \textbf{Se mantiene.} Quedan nueve exclusiones con los códigos del SD2 y un análisis del riesgo que protege cada una. & SD3, §3.2.2, p. 10; Anexos, Anexo 3.D, p. 134. \\
\LFXfiletefila
% ID: SD3-12
\mbox{SD3-12} & No cita la objeción de la gerenta general, no analiza dependencias entre épicas y el hito de septiembre de 2026 no figura. & \textbf{Documentado.} Se acoge la objeción de la gerenta: las rutas se adelantan a la Etapa 1 por la jubilación del planificador. Las dependencias se analizan por grupo de familias. Se trata el vencimiento de la suspensión del proveedor de lácteos en septiembre de 2026, junto con la consulta V-13, ahora con cita al caso (cap. 13.1 y 13.2, p. 22). & SD3, §3.1.2, p. 7–8; §3.2.1, p. 8–10. \\
\LFXfiletefila
% ID: SD3-13
\mbox{SD3-13} & Contradicción con el Art. 17°: se decía «Marcha blanca de la Etapa 2… Meses 13–20». & \textbf{Documentado.} La Tabla 3.1 fija la Etapa 2 con desarrollo en los meses 13–18, marcha blanca en los meses 19–20 y producción en el mes 21. También se analiza la compatibilidad del calendario con las prohibiciones de septiembre y diciembre (S-17 y V-17). & SD3, Tabla 3.1 y §3.1.2, p. 7–8. \\
\LFXfiletefila
% ID: SD3-14
\mbox{SD3-14} & Las etapas no cuadran: OTIF y fill rate se asignaban a la Etapa 2 con RF en la Etapa 1, y el EDI de la Etapa 2 se habilitaba en los meses 10–16. & \textbf{Documentado.} OTIF y fill rate quedan en la Etapa 1 (RF-11.03 y RF-11.04, E1; D-01 y D-02; M10). La Etapa 2 conserva solo el costo de servir. No quedan menciones a «meses 10–16». El canal moderno entra en producción en el mes 21 (RNF-12.01). & SD3, §3.2.1, p. 9; Tabla 3.4; Anexos, Tabla 3.A.1 (RF-11.03 y RF-11.04) y Tabla 3.A.9 (D-01 y D-02). \\
\LFXfiletefila
% ID: SD3-15
\mbox{SD3-15} & Exclusiones y restricciones infladas con relleno licitatorio (EX-12 a EX-14, R-14 a R-27). & \textbf{Documentado.} Quedan nueve exclusiones (E-01 a E-09) y doce restricciones (R-01 a R-12), todas del caso. & Anexos, Tablas 3.A.7 y 3.A.8, p. 134–136. \\
\LFXfiletefila
% ID: SD3-16
\mbox{SD3-16} & Los supuestos no tienen identificador, fundamento ni validación; se cita un «S-41» inexistente. & \textbf{Documentado.} Hay 44 supuestos (S-01 a S-44) con columnas de supuesto, fundamento, impacto si resulta falso, validación y origen. S-41 existe ahora: computadores de 184 personas de administración. & SD3, §3.2.2, p. 10–11; Anexos, Tabla 3.A.6, p. 117–134. \\
\LFXfiletefila
% ID: SD3-17
\mbox{SD3-17} & No hay visión de arquitectura empresarial; interoperabilidad, escalabilidad, seguridad y performance aparecen dispersas. & \textbf{Documentado.} §3.4.2 da la vista por capas y ahora cierra con una síntesis de cuatro atributos: interoperabilidad (ERP como emisor único, cola durable de §3.3.3), escalabilidad (RNF-19.01, 19.02, 19.04; RT-02.12), desempeño (RNF-05.01, 03.01, 03.02, 11.02; primer cuello de botella en la emisión de guías del ERP) y seguridad (RNF-14.01, 14.02, 14.04, 14.06, 14.07, 15.01, 15.02). Remite el detalle al SD4, secciones 4.1.3.7 y 4.2.6. & SD3, §3.3.3, p. 18; §3.4.2, p. 26–27. \\
\LFXfiletefila
% ID: SD3-18
\mbox{SD3-18} & Crítico: el Formulario T-12 no fue entregado. & \textbf{Documentado.} Es un formulario independiente, con 645 filas: Parte A (RF y RNF) y Parte B (RT), incluida la columna EDT. & Formulario T-12. \\
\LFXfiletefila
% ID: SD3-19
\mbox{SD3-19} & El catálogo funcional es una lista de nombres, sin descripción verificable, actor, precondición ni resultado esperado. & \textbf{Documentado.} Cada RF tiene actor, descripción, precondición, resultado esperado, prioridad y etapa, y origen, según el 17.1 del caso (p. 31). & SD3, §3.2.3, p. 11; Anexos, Anexo 3.A, p. 6–88. \\
\LFXfiletefila
% ID: SD3-20
\mbox{SD3-20} & La prioridad no discrimina: 80 de 90 RF son «Alta» o «Muy Alta». & \textbf{Documentado.} Hay una regla explícita: la prioridad no supera la de la familia del SD2 y «Crítica» queda solo para resultados de aceptación, restricciones e hitos. Se mantienen 56 críticos, 163 altos y 43 medios (84 \% Crítica/Alta), y §3.2.3 justifica esa concentración: el catálogo describe el flujo operativo que el Caso §13.1 pone primero; costo de servir y autoatención (F-12, F-14) son Media; «Media» no es optativo (T-13, T-17). & SD3, §3.2.3, p. 11–12. \\
\LFXfiletefila
% ID: SD3-21
\mbox{SD3-21} & La numeración tiene saltos: parece el extracto de un catálogo mayor. & \textbf{Documentado.} En el T-12 ninguna épica RF tiene saltos de numeración. Las nueve materias absorbidas o alias se declaran como tales. & Formulario T-12, Parte A; Anexos, Tablas 3.A.3a y 3.A.5a. \\
\LFXfiletefila
% ID: SD3-22
\mbox{SD3-22} & El catálogo no funcional no tiene método de verificación. & \textbf{Documentado.} Hay una columna «Verificación prevista» en los RNF y la prueba asociada se registra en el T-13 §8.2. & Anexos, Anexo 3.B, p. 89–116. \\
\LFXfiletefila
% ID: SD3-23
\mbox{SD3-23} & Se citan requerimientos inexistentes: RF-02.05d/e y RF-15.05. & \textbf{Documentado.} No queda ninguna mención a RF-02.05d/e. RF-15.05 existe ahora en el catálogo (aprovisionamiento de cuentas) y en el T-12. & Anexos, Tabla 3.A.3 (RF-15.05); Formulario T-12. \\
\LFXfiletefila
% ID: SD3-24
\mbox{SD3-24} & No existe la sección de implementación: ni metodología, ni pipeline, ni gestión de configuración, ni ramas. & \textbf{Documentado.} §3.4.3 cubre la base compartida, tres frentes por etapa, el ciclo de cada incremento (rama versionada con configuración, revisión de pares, CI con cobertura bloqueante, promoción DEV → QA → PREPROD y pruebas de perfil operacional), los criterios que impiden promover y el orden de construcción. El umbral de cobertura quedó alineado con D-12: 70 \% de lógica de negocio y 80 \% unitario global, ambos bloqueantes. & SD3, §3.4.3, p. 27–28. \\
\LFXfiletefila
% ID: SD3-25
\mbox{SD3-25} & No existe la sección de implantación: las oleadas y la migración eran filas de tabla. & \textbf{Documentado.} §3.4.4 desarrolla las siete condiciones del caso (cap. 13.3, p. 23), la secuencia de olas, el criterio de avance, la estabilización con dotación calculada (2 + 3 + 7 personas), la adopción por perfil y la migración por sitio con conciliación. & SD3, §3.4.4, p. 28–30. \\
\LFXfiletefila
% ID: SD3-26
\mbox{SD3-26} & Se habla de «plan de reversión» sin describirlo ni dar tiempo de rollback para la ventana de 05:30 a 07:00. & \textbf{Documentado.} La reversión tiene responsable que la autoriza, causales y un objetivo de 40 minutos: decisión a las 04:45 y término a las 05:25 (T-18 §6.4), ensayado en PREPROD antes de cada corte. & SD3, §3.4.4, p. 29. \\
\LFXfiletefila
% ID: SD3-27
\mbox{SD3-27} & No hay convivencia con el papel durante la marcha blanca (13.3 y 17.6 del caso). & \textbf{Documentado.} La hoja de picking y los documentos en papel quedan como consulta y apoyo durante la ola, y se retiran tras cuatro semanas cumpliendo el criterio de avance. & SD3, §3.4.4, p. 28–29. \\
\LFXfiletefila
% ID: SD3-28
\mbox{SD3-28} & No existe la sección de operación: «SRE» aparece una sola vez y no hay objetivos de servicio, runbooks ni automatización. & \textbf{Documentado.} §3.4.5 tiene la Tabla 3.5 de objetivos, monitoreo, escalamiento y peaks, y ahora declara el enfoque SRE: los objetivos deciden prioridades (estabilidad antes que cambios nuevos al acercarse al límite), infraestructura como código, alertas preventivas con escalamiento automático, escalamiento automático en peaks (RNF-19.02), bloqueo automático en la cadena de integración, libro de operación con automatización progresiva (RNF-16.02) y revisión de causa tras cada incidente crítico o alto. & SD3, §3.4.5, p. 30–32. \\
\LFXfiletefila
% ID: SD3-29
\mbox{SD3-29} & La mesa de ayuda funciona de 08:00 a 20:00, pero el caso exige de 04:00 a 22:00 de lunes a sábado (RT-21.06). & \textbf{Documentado.} El horario queda de 04:00 a 22:00 de lunes a sábado, 24×7 en septiembre y diciembre, con críticos y altos 24×7×365, ahora con cita al caso (cap. 15, RT-21.06, p. 27). & SD3, §3.4.5, p. 30 y 32. \\
\LFXfiletefila
% ID: SD3-30
\mbox{SD3-30} & OK: RTO ≤ 4 h y RPO ≤ 15 min coherentes. & \textbf{Se mantiene.} & SD3, Tabla 3.5, p. 31. \\
\LFXfiletefila
% ID: SD3-31
\mbox{SD3-31} & OK: la tabla de grupos de interés tiene cuadrante, momento, responsable e indicador. & \textbf{Se mantiene y se amplía.} Se agregó una respuesta ante la resistencia de cada actor. & Anexos, Tabla 3.A.12, p. 152–157. \\
\LFXfiletefila
% ID: SD3-32
\mbox{SD3-32} & Faltan la gerenta general, el comité ejecutivo, el gerente comercial, los almaceneros, la autoridad sanitaria y los peonetas. & \textbf{Documentado.} Están los 19 actores del SD2, con el comité presidido por la gerenta general. La Tabla 3.A.12a separa a los 15 actores del sistema de los interesados sin acceso. & SD3, §3.4.6, p. 32–33; Anexos, Tablas 3.A.12 y 3.A.12a, p. 152–160. \\
\LFXfiletefila
% ID: SD3-33
\mbox{SD3-33} & El planificador quedaba en «C — Mantener informado», aunque el criterio N.° 16 depende de él. & \textbf{Documentado.} El planificador figura con influencia Alta e interés Muy Alto y participa desde el mes 1 en talleres de reglas, no solo como informado. & SD3, §3.4.6, p. 32–33; Anexos, Tabla 3.A.12. \\
\LFXfiletefila
% ID: SD3-34
\mbox{SD3-34} & OK: los 16 criterios tienen línea base, meta, momento y medición, pero el N.° 16 no define cómo se verifica. & \textbf{Documentado.} Para R18-16 se compara el OTIF de dos semanas de planificación sin el planificador contra dos semanas comparables previas, durante la marcha blanca de la Etapa 1. & Anexos, Tabla 3.A.13, p. 160–163. \\
\LFXfiletefila
% ID: SD3-35
\mbox{SD3-35} & OK: las 16 decisiones del 16.1 están tomadas, pero enterradas en 65 páginas de tabla. & \textbf{Documentado.} Las 16 decisiones se registran como supuestos S-01 a S-16, alineados con el SD2, y el cuerpo las resume en un punteo (§3.2.2). Las decisiones de diseño se redujeron a 12 (D-01 a D-12). Las reglas sensibles se explican en §3.2.4. El informe valoró además el reemplazo del WMS de 2013 por estrangulamiento (D14), que se mantiene en el SD4 (ADR-08), y el registro de vacíos y consultas que detectó errores de las Bases, que se conserva como Anexo 3.H. & SD3, §3.2.2 y §3.2.4, p. 10–13; Anexos, Tablas 3.A.6 y 3.A.9; Anexos, Anexo 3.H, p. 145; SD4, §4.1.3.5 (ADR-08). \\
\LFXfiletefila
% ID: SD3-36
\mbox{SD3-36} & La excursión térmica es contradictoria: «lo decide el conductor» frente a «la adopta la autoridad de calidad»; el SD4.1 dice que bloquea el sistema. & \textbf{Documentado.} Hay una regla única (RNG-04 y S-04): advertencia ante una excursión menor y bloqueo preventivo del sistema ante una crítica y sostenida, con la disposición final siempre a cargo de Calidad. El conductor no decide el destino sanitario. Mientras Calidad no apruebe los parámetros de un tipo de producto, toda lectura fuera del rango de su ficha se trata como crítica; esa regla ya no se presenta como «pendiente de validación». & SD3, §3.2.2, p. 10–11; §3.2.4, p. 12–13; §3.4.1, p. 18–19; Anexos, RNG-04 y S-04. \\
\LFXfiletefila
% ID: SD3-37
\mbox{SD3-37} & Precio entre pedido y despacho contradictorio: el supuesto decía «precio ofrecido» y D9/RNG-08 decían «vigente al despacho». & \textbf{Documentado.} Rige el precio acordado al tomar el pedido, en S-09 y RNG-08, igual que en el SD2. Una diferencia con el ERP bloquea la facturación hasta conciliarse. & SD3, §3.2.2, p. 10–11; Anexos, RNG-08 y S-09. \\
\LFXfiletefila
% ID: SD3-38
\mbox{SD3-38} & La promesa de 24 h se da por verificable sin decisión, requerimiento ni criterio. & \textbf{Documentado.} Queda en RNG-15, D-03 y S-21: 24 h para pedidos urbanos antes de las 14:00 y 48 h para el resto. La preventa muestra la fecha prometida y el OTIF se mide contra ella (RNG-09). & SD3, §3.2.2 y §3.2.4, p. 10–13. \\
\LFXfiletefila
% ID: SD3-39
\mbox{SD3-39} & Crítico: D45 incluía un precio («113 mil dólares…»), en contradicción con EX-09. & \textbf{Documentado.} No queda ningún precio en ninguna moneda en el cuerpo, los anexos ni el T-12. La compra del hardware por el CLIENTE se mantiene (E-09). & SD3, §3.2.2, p. 10; Anexos, Tabla 3.A.7. \\
\LFXfiletefila
% ID: SD3-40
\mbox{SD3-40} & Se citan nombres de archivo como fuente («Caso\_02\_Logistica.md», etc.). & \textbf{Documentado.} No queda ningún nombre de archivo ni la denominación «Caso 02». Las Bases se citan por título abreviado con capítulo o artículo y página (\emph{Bases Administrativas}, \emph{Bases Técnicas Transversales} y \emph{Caso: Logística}). & SD3, Referencias, p. 33; Anexos, Referencias, p. 167. \\
\LFXfiletefila
% ID: SD3-41
\mbox{SD3-41} & Quedaron notas de corrección de una versión anterior («La decisión previa… se reformula», «no estaba respondida»…). & \textbf{Documentado.} No quedan esas notas. Se quitaron además las fechas de edición de las tablas de IA del cuerpo y de los anexos, la nota «esta actualización» de los anexos y la frase de proceso de las Referencias del cuerpo. & SD3, Tabla 3.6, p. 34–35; Anexos, Tabla 3.A.15, p. 167. \\
\LFXfiletefila
% ID: SD3-42
\mbox{SD3-42} & Dos voces distintas (D1–D40 frente a D41–D60). & \textbf{Documentado.} El registro de 60 decisiones se reescribió como D-01 a D-12 en una sola voz; las decisiones D-08, D-09, D-11 y D-12, que seguían en estilo telegráfico, se redactaron como oraciones completas igual que las demás. El cuerpo no usa primera persona y nombra a LafroX y al CLIENTE de forma uniforme. & Anexos, Tabla 3.A.9, p. 139–140. \\
\LFXfiletefila
% ID: SD3-43
\mbox{SD3-43} & Se citaban decisiones antes de que existieran, había códigos sin glosario (BTC, BTT, R18, RNG, HHT) y un párrafo de excursión pegado en la decisión de la cámara a −22 °C. & \textbf{Documentado.} No quedan citas a «Decisión 20» ni «Decisión 29». El Anexo 3.K es un glosario de siglas y códigos. Las decisiones vigentes no mezclan materias. & Anexos, Anexo 3.K, p. 164–166; Tabla 3.A.9. \\
\LFXfiletefila
% ID: SD3-44
\mbox{SD3-44} & Erratas de ensamblaje: «ventana de 320 minutos» (p. 3) y «160 camiones que no son trabajadores» (p. 66). & \textbf{Documentado.} Ninguna de las dos frases aparece en el cuerpo, los anexos ni el T-12; tampoco la errata «próposito». & SD3, anexos y T-12. \\
\LFXfiletefila
% ID: SD3-45
\mbox{SD3-45} & Títulos huérfanos al pie de página y páginas semivacías (21 en el SD3). & \textbf{Documentado.} Se revisaron página por página el cuerpo, los anexos y el T-12. Ningún título de sección queda al pie de página y las páginas semivacías se redujeron de 21 a menos de diez, debidas a figuras y tablas en hoja horizontal, a filas de tabla que no se dividen y a páginas de índice. & SD3, anexos y T-12. \\
\end{longtable}
}
\Needspace{18\baselineskip}
\subsection{Subdocumento 4}
{\linespread{1}\LFXcuerpotabla
\setlength{\parskip}{0pt}
\renewcommand{\arraystretch}{1.10}
\setlength{\tabcolsep}{4pt}
\begin{longtable}{F{0.10} F{0.27} F{0.43} F{0.20}}
\caption{Observaciones y respuestas --- subdocumento 4}\\[2pt]
\LFXtopruleoriginal \cab{N.º} & \cab{Observación del Informe 1} & \cab{Respuesta} & \cab{Sección modificada} \\ \midrule
\endfirsthead
\multicolumn{4}{@{}l@{}}{\fontsize{9pt}{11pt}\selectfont\color{lafroxGris}Tabla \thetable\ --- Subdocumento 4 (continuación)}\\[3pt]
\LFXtopruleoriginal \cab{N.º} & \cab{Observación del Informe 1} & \cab{Respuesta} & \cab{Sección modificada} \\ \midrule
\endhead
\midrule
\multicolumn{4}{@{}r@{}}{\fontsize{9pt}{11pt}\selectfont\color{lafroxGris}Continúa en el folio siguiente.}\\
\endfoot
\bottomrule
\endlastfoot
% ID: SD4-01
\mbox{SD4-01} & \textbf{EL-01} — Texto y diagramas desconectados, pp. 9–10 & Las figuras de las capas se citan antes de aparecer y se explican inmediatamente después, vinculando el recorrido gráfico con el texto. \textbf{Estado y cierre:} Documentado; comprobar legibilidad final. & SD4, §4.1.3; figuras 1–12. \\
\LFXfiletefila
% ID: SD4-02
\mbox{SD4-02} & \textbf{EL-02} — Trece recorridos por actor redundantes y sin explicación, p. 9 & Los trece recorridos por actor se sustituyen por una vista general, ocho vistas por capa, secuencias y una matriz de actores. \textbf{Estado y cierre:} Sustitución documentada; cotejar permiso local de Calidad en figuras. & SD4, §4.1.3; Anexo 4-N, matriz de actores. \\
\LFXfiletefila
% ID: SD4-03
\mbox{SD4-03} & \textbf{EL-03} — Actores declarados distintos de los dibujados, p. 9 & Se identifican quince actores y se describen sus perfiles y ámbitos de autorización. La coincidencia con las figuras todavía requiere cotejo gráfico. \textbf{Estado y cierre:} Cotejo gráfico final pendiente. & SD4, §§4.1.3.1 y 4.1.18; Anexo 4-N. \\
\LFXfiletefila
% ID: SD4-04
\mbox{SD4-04} & \textbf{EL-04} — Matriz requisito → módulo → capa ausente, p. 9 & El T-12 identifica los requisitos individualmente; 4-E agrupa familias y 4-N relaciona módulos, 37 componentes transversales y realización física. La Tabla 6 y el T-11 completan el emplazamiento y los recursos. \textbf{Estado y cierre:} Cadena documentada; cotejo individual de suficiencia en curso y precisión de sus referencias externas por cotejar. & Formulario T-12; anexos 4-E y 4-N, Tabla A.14; SD4, §4.2.2, Tabla 6; Formulario T-11. \\
\LFXfiletefila
% ID: SD4-05
\mbox{SD4-05} & \textbf{EL-05} — Cinco ambientes solo nombrados, p. 9 & Los cinco ambientes tienen desarrollo por separado y figuras propias. La topología AWS identifica las regiones de producción y recuperación. \textbf{Estado y cierre:} Documentado por ambiente. & SD4, §§4.1.9 y 4.2.4.1; figuras 19–23 y Figura 18. \\
\LFXfiletefila
% ID: SD4-06
\mbox{SD4-06} & \textbf{EL-06} — SOLID y patrones sin aplicación, pp. 9–10 & Se desarrolla la aplicación de SOLID mediante contratos de negocio y límites entre módulos, junto con los patrones de diseño y continuidad. \textbf{Estado y cierre:} Documentado; pruebas de dependencia/contrato por ejecutar. & SD4, §§4.1.10 y 4.1.10.1. \\
\LFXfiletefila
% ID: SD4-07
\mbox{SD4-07} & \textbf{EL-07} — Seis vistas declaradas sin vista de procesos, p. 10 & Se establece la correspondencia de las cinco vistas exigidas y se incorporan secuencias de procesos. \textbf{Estado y cierre:} Corregido documentalmente. & SD4, §4.1.14; figuras 8–9; Anexo 4-S. \\
\LFXfiletefila
% ID: SD4-08
\mbox{SD4-08} & \textbf{EL-08} — Estilo modular pertinente al caso, p. 10 & Se mantiene el estilo modular valorado por el evaluador; Laravel/PHP conserva límites entre los módulos y su elección queda justificada en ADR-01. \textbf{Estado y cierre:} Observación positiva conservada; framework actualizado. & SD4, §4.1.13; Anexo 4-O, ADR-01. \\
\LFXfiletefila
% ID: SD4-09
\mbox{SD4-09} & \textbf{EL-09} — Comparación insuficiente y ADR fuera de lógica, p. 10 & Se comparan alternativas arquitectónicas en el cuerpo y se conserva un registro único de 22 ADR en el anexo lógico, sin repetir sus fichas en el cuerpo. \textbf{Estado y cierre:} Documentado, sin repetir fichas en cuerpo. & SD4, §4.1.13, Tabla 4; Anexo 4-O. \\
\LFXfiletefila
% ID: SD4-10
\mbox{SD4-10} & \textbf{EL-10} — Doce módulos anunciados sin inventario, p. 10 & Los doce módulos cuentan con inventario de nombre, responsabilidad, contratos, actor y etapa. \textbf{Estado y cierre:} Documentado. & Anexo 4-D. \\
\LFXfiletefila
% ID: SD4-11
\mbox{SD4-11} & \textbf{EL-11} — Límites, interfaces, orquestación y coreografía ausentes, p. 10 & Se desarrollan límites de contexto, interfaces internas y externas, y la relación entre orquestación y coreografía mediante contratos y secuencias. \textbf{Estado y cierre:} Documentado. & SD4, §4.1.3.5; figuras 8–9; anexos 4-F, 4-G y 4-H. \\
\LFXfiletefila
% ID: SD4-12
\mbox{SD4-12} & \textbf{EL-12} — Marco TOGAF solo citado, p. 10 & La descripción aplica ISO 42010 mediante interesados, preocupaciones y correspondencias entre vistas. No se presenta la mención de TOGAF como evidencia de aplicación. \textbf{Estado y cierre:} Se usa el marco desarrollado; no se afirma aplicación de TOGAF por su sola mención. & Anexo 4-S. \\
\LFXfiletefila
% ID: SD4-13
\mbox{SD4-13} & \textbf{EL-13} — Modelo de dominio solo como lista, p. 10 & Se incorpora un modelo conceptual del dominio y se relacionan sus productores con los eventos canónicos. El esquema relacional se remite al SD5. \textbf{Estado y cierre:} Documentado; esquema relacional pertenece a SD5. & SD4, §4.1.5, Figura 13; Anexo 4-A. \\
\LFXfiletefila
% ID: SD4-14
\mbox{SD4-14} & \textbf{EL-14} — Caché de 8 h incompatible con 24 h y relevos, p. 10 & El diseño usa caché de 24 horas, manifiesto de 26 horas, credenciales personales y verificador local. AL-OFF-01 incluye relevos a las 8 y 16 horas durante el corte. \textbf{Estado y cierre:} Diseño documentado; probar vigencias, permisos, reinicio y reconciliación; una copia vigente no renueva credenciales ni autoriza terceros. & SD4, §§4.1.3.7 y 4.1.15; anexos 4-J y 4-M, AL-OFF-01. \\
\LFXfiletefila
% ID: SD4-15
\mbox{SD4-15} & \textbf{EL-15} — Gateway solo remoto para HHT, p. 10 & La puerta de API local A-01 se ejecuta en VM-01, VM-C01 y E-01, de modo que el HHT no depende de API Gateway durante un corte de WAN. \textbf{Estado y cierre:} Documentado; no depende de API Gateway durante corte. & SD4, §§4.1.3.3 y 4.2.2; Figura 5; Anexo 4-N. \\
\LFXfiletefila
% ID: SD4-16
\mbox{SD4-16} & \textbf{EL-16} — Bloqueo térmico decidido por conductor frente a automático, p. 10 & M9 bloquea preventivamente una excursión crítica y sostenida; la disposición sanitaria corresponde a Calidad con evidencia y permisos. \textbf{Estado y cierre:} Bloqueo documentado; cotejar etiquetas gráficas y recorrido local autorizado de Calidad antes de cerrar la autonomía de disposición. & SD4, §4.1.4.2; anexos 4-J, 4-M y 4-N. \\
\LFXfiletefila
% ID: SD4-17
\mbox{SD4-17} & \textbf{EL-17} — Etapas discordantes con SD3, p. 10 & Se declaran las etapas por módulo, la implantación progresiva y la remisión al calendario. La consistencia con el SD3 vigente permanece por comprobar. \textbf{Estado y cierre:} Contrastar contra SD3 vigente sin editarlo. & SD4, §4.1.8; Anexo 4-D; contraste externo con SD3 y calendario. \\
\LFXfiletefila
% ID: SD4-18
\mbox{SD4-18} & \textbf{EL-18} — Portales externos tratados como intranet/VPN, p. 10 & Se distinguen entradas públicas, privadas y locales. Verified Access se reserva a oficina y consolas, sin exigir VPN para los portales externos públicos. \textbf{Estado y cierre:} Documentado; Verified Access reservado a oficina/consolas. & SD4, §§4.1.3 y 4.2.5.2. \\
\LFXfiletefila
% ID: SD4-19
\mbox{SD4-19} & \textbf{EL-19} — Portada promete física pero solo incluye lógica, p. 10 & El cuerpo integrado incluye arquitectura lógica, física y centros de datos bajo la estructura 4.1–4.3. \textbf{Estado y cierre:} Corregido en estructura. & SD4, portada y §§4.1–4.3; ensamblador integrado. \\
\LFXfiletefila
% ID: SD4-20
\mbox{SD4-20} & \textbf{EL-20} — Códigos y versiones sin documento entregado, p. 10 & Se identifica el catálogo físico A/B/C/D/E/F/N, el inventario lógico y los contratos internos y externos. Queda comprobar que las imágenes no conserven códigos históricos. \textbf{Estado y cierre:} Verificar etiquetas residuales en imágenes; no mezclar códigos históricos. & SD4, §4.2.2; anexos 4-N, 4-G y 4-H; Formulario T-11. \\
\LFXfiletefila
% ID: SD4-21
\mbox{SD4-21} & \textbf{EL-21} — Plantilla distinta e ilegibilidad, p. 10 & Los tres entregables usan plantilla común. El cuerpo conserva 26 tablas técnicas breves y su declaración de IA; la red se presenta mediante tres activos gráficos actuales. \textbf{Estado y cierre:} Formato documentado; el registro histórico de rótulos próximos a 5 puntos no describe por sí solo los nuevos activos de red. Revisar su impresión final. & SD4, cuerpo y figuras; anexos y T-11; registro AUDITORIA\_TABLAS\_SD4.md. \\
\LFXfiletefila
% ID: SD4-22
\mbox{SD4-22} & \textbf{EF-01} — Documento duplicado y sin portada; entrega con índice, folio y firma, pp. 11 y 13 & Se organiza un cuerpo integrado, con anexos y T-11 como documentos independientes, evitando la duplicación estructural de la entrega preliminar. \textbf{Estado y cierre:} Estructura corregida; verificar índice paginado, folios y firma de los archivos finales. & SD4, portada, índice, folios y firma; documentos de anexos y Formulario T-11. \\
\LFXfiletefila
% ID: SD4-23
\mbox{SD4-23} & \textbf{EF-02} — Comentarios de IA, versiones y auditorías dentro de la oferta, p. 11 & El barrido previo separó las notas de trabajo de la oferta. Se conserva la declaración obligatoria de IA con el revisor identificado en el estado actual. El PDF vigente no contiene «Cambios v01 → v02», «Si se prefiere (B)», «hallazgo D1», «3 brokers», «En texto, la vista» ni referencias a versiones como «v06» o «v6.2». \textbf{Estado y cierre:} La declaración actual identifica a Álvaro Catalán en 132 filas; comprobar ausencia de notas de proceso en la revisión final, sin retirar la declaración obligatoria. & SD4, declaración de uso de IA; registros de revisión externos al entregable. \\
\LFXfiletefila
% ID: SD4-24
\mbox{SD4-24} & \textbf{EF-03} — Tabla de emplazamiento inexistente/incompleta, p. 11 & El emplazamiento se relaciona con los doce módulos y las 37 capacidades transversales. La matriz de correspondencia, el inventario y el T-11 identifican su realización física. \textbf{Estado y cierre:} Doce módulos y 37 transversales inventariados; cotejo individual necesario para confirmar realización física y cobertura de requisitos. & SD4, §4.2.2, Tabla 6; Anexo 4-N, Tabla A.14; Anexo 4-W, Tabla A.35; T-11. \\
\LFXfiletefila
% ID: SD4-25
\mbox{SD4-25} & \textbf{EF-04} — Figura física ilegible y sin explicación, p. 11 & La vista física general se acompaña de explicaciones y vistas de sitios, emplazamiento, nube y recuperación. Las tres vistas de red detallan la conexión local y regional. \textbf{Estado y cierre:} Explicación presente; cotejar correspondencia con recuperación y legibilidad actual. & SD4, §§4.2 y 4.3; Figura 14, figuras 15–18 y 23–26. \\
\LFXfiletefila
% ID: SD4-26
\mbox{SD4-26} & \textbf{EF-05} — Descripciones verbales sustituyen dibujos, pp. 12–13 & Se incorporan dibujos de ambientes, nube, sitios, racks y recinto. Las tres vistas de red forman parte del desarrollo y tienen explicación propia. \textbf{Estado y cierre:} Incorporación documental corregida; queda revisión de correspondencia gráfica y legibilidad. El borrador de propuestas no es el activo entregado. & SD4, §4.2.4.2, figuras 24–26; §4.3.1.4, figuras 28–29; vistas de ambientes y sitios. \\
\LFXfiletefila
% ID: SD4-27
\mbox{SD4-27} & \textbf{EF-06} — BI ocupa catorce páginas ajenas al ítem, p. 12 & El SD4 conserva la función lógica y el emplazamiento de M10. El modelo analítico detallado se remite al SD5. \textbf{Estado y cierre:} Mantener solo justificación y dependencias arquitectónicas. & SD4, §§4.1.4.7 y 4.2.2; remisión a SD5. \\
\LFXfiletefila
% ID: SD4-28
\mbox{SD4-28} & \textbf{EF-07} — Implementos pertinentes, p. 12 & Se mantiene el contenido valorado positivamente: el T-11 identifica modelos de implementos, cantidades y finalidad. \textbf{Estado y cierre:} Observación positiva conservada. & Formulario T-11. \\
\LFXfiletefila
% ID: SD4-29
\mbox{SD4-29} & \textbf{EF-08} — Cantidades sin número y ciclo de vida/caídas omitidos, p. 12 & Se explicitan cantidades y criterios de aceptación por tipología, junto con los cálculos del parque y el control de actualizaciones. La aceptación física debe verificarse por ensayo. \textbf{Estado y cierre:} Comprobar especificación de aceptación por tipo; cálculos no equivalen a ensayo. & SD4, §4.2.1.2; Anexo 4-W.3; Anexo 4-P; Formulario T-11. \\
\LFXfiletefila
% ID: SD4-30
\mbox{SD4-30} & \textbf{EF-09} — Exceso de equipos congelados/POS/impresoras, p. 12 & El parque se deriva de la dotación por turno y sus reservas: 207 terminales al inicio y 235 en el año 3, incluidos repuestos. La concurrencia nocturna inicial es 186. \textbf{Estado y cierre:} Cifras derivadas; 186 es concurrencia nocturna inicial. Precisar el perfil temporal de la campaña de Preproducción que menciona 207. & SD4, §§4.2.1 y 4.1.9; Anexo 4-W.3; Formulario T-11. \\
\LFXfiletefila
% ID: SD4-31
\mbox{SD4-31} & \textbf{EF-10} — Frío solo en 18 camiones y alertas al regreso, p. 12 & Se declaran 28 termógrafos: 18 en camiones propios y 10 en camiones externos. La lectura BLE y la alerta local se producen durante la ruta. \textbf{Estado y cierre:} Corregido en diseño; ensayo de lectura/alerta por ejecutar. & SD4, §4.1.4.2; Anexo 4-W.4; Formulario T-11, B-03. \\
\LFXfiletefila
% ID: SD4-32
\mbox{SD4-32} & \textbf{EF-11} — Estudio de cobertura inalámbrica insuficiente, p. 12 & Se incluye el estudio de cobertura y la aceptación de las interfaces en cámaras de frío. La propuesta no presenta la estimación como cobertura ya medida. \textbf{Estado y cierre:} Estudio y medición necesarios; no inventar cobertura medida. & SD4, §4.2.1.2; criterios de aceptación de interfaces y cámaras. \\
\LFXfiletefila
% ID: SD4-33
\mbox{SD4-33} & \textbf{EF-12} — Sala secundaria sobredimensionada y Tier incoherente, p. 12 & El dimensionamiento usa el equipamiento real del recinto y distingue disponibilidad por componente de disponibilidad de extremo a extremo. \textbf{Estado y cierre:} Documentado; ninguna categoría Tier sustituye demostrar disponibilidad. & SD4, §4.3.1.4 y desarrollo de disponibilidad. \\
\LFXfiletefila
% ID: SD4-34
\mbox{SD4-34} & \textbf{EF-13} — Sin cálculo eléctrico/térmico, plano ni racks, p. 12 & Se desarrollan cargas eléctricas y térmicas, memoria de cálculo, distribución de racks y plano del recinto. \textbf{Estado y cierre:} Evidencias presentes; verificar correspondencia gráfica final y actualizar referencias externas del T-12. & SD4, §4.3.1.4, Tablas 23–24 y figuras 28–29; Anexo 4-W.8. \\
\LFXfiletefila
% ID: SD4-35
\mbox{SD4-35} & \textbf{EF-14} — Gabinetes de borde sin energía/clima/acceso, p. 12 & Los gabinetes de borde cuentan con representación gráfica y definición de energía, climatización, acceso y carga/UPS por tipología. \textbf{Estado y cierre:} Documentado por tipología. & SD4, §4.2.1.2, figuras 15–16; Anexo 4-W, Tabla A.39; Formulario T-11. \\
\LFXfiletefila
% ID: SD4-36
\mbox{SD4-36} & \textbf{EF-15} — Aprobación de DR podría llevar RTO a 24–72 h, p. 12 & Se declaran las regiones, RTO de hasta 4 horas y aprobación de residencia como hito contractual. La consulta sanitaria regional de solo lectura tiene un presupuesto de 75 minutos, sujeto a prueba. \textbf{Estado y cierre:} Diseño documentado; medir restitución por servicio, conectividad de Talca y consulta sanitaria; residencia y brecha extrema de RPO requieren decisión contractual real. & SD4, §4.3.2; Anexo 4-M, AL-DR-01; Formulario T-11. \\
\LFXfiletefila
% ID: SD4-37
\mbox{SD4-37} & \textbf{EF-16} — T-11 invocado pero no adjunto, p. 12 & El T-11 existe como formulario independiente. \textbf{Estado y cierre:} Presente; comprobar su inclusión en el paquete de entrega. & Formulario T-11, archivo independiente. \\
\LFXfiletefila
% ID: SD4-38
\mbox{SD4-38} & \textbf{EF-17} — Zero Trust y SPOF pertinentes, p. 12 & Se mantienen Zero Trust, el análisis de puntos únicos de falla y el respaldo 3-2-1-1-0, con remisión a su desarrollo principal. \textbf{Estado y cierre:} Contenido conservado; remitir a sección principal sin repetir controles. & SD4, §§4.1.3.7 y 4.2.5; anexos 4-Q y 4-R; respaldo y continuidad. \\
\LFXfiletefila
% ID: SD4-39
\mbox{SD4-39} & \textbf{EF-18} — Identidad, gateway local y funciones offline sin resolver, p. 13 & La identidad local se explica en 4.1.3.7 y la puerta local en 4.1.3.3. Se distinguen funciones disponibles y bloqueadas durante un corte: captura local, relevos, guía preemitida y autorizaciones externas condicionadas. \textbf{Estado y cierre:} Funciones y límites documentados; probar continuidad completa, no solo captura; completar cotejo del acceso local de Calidad. & SD4, §§4.1.3.3, 4.1.3.7 y 4.1.15–4.1.16; anexos 4-J, 4-K y 4-M. \\
\LFXfiletefila
% ID: SD4-40
\mbox{SD4-40} & \textbf{EF-19} — Ceph sobre RAID y capacidad contradictoria, p. 13 & Se define Ceph con seis OSD sin RAID, size=3/min\_size=2, redundancia N+1 y cálculo de capacidad. ADR-10 contrasta la demanda de 136 IOPS actuales y 359 a 3× con la capacidad del NVMe, que no es criterio de selección ni justifica equipos adicionales. \textbf{Estado y cierre:} Diseño documentado; fallo y reconstrucción por probar. & SD4, arquitectura física de Talca; Anexo 4-W.8; Anexo 4-O, ADR-10. \\
\LFXfiletefila
% ID: SD4-41
\mbox{SD4-41} & \textbf{EF-20} — DTE diferido contradice único emisor/guía previa, p. 13 & El ERP conserva la emisión previa de la guía. Un cambio de carga invalida la preemisión y requiere otra guía válida; AL-DTE-01 define el ensayo. \textbf{Estado y cierre:} Documentado; 96 salidas requieren prueba completa y documentos válidos. & SD4, §4.1.17; anexos 4-M y 4-U, AL-DTE-01. \\
\LFXfiletefila
% ID: SD4-42
\mbox{SD4-42} & \textbf{EF-21} — Dotación incompatible con cuatro TI, p. 13 & El dimensionamiento incorpora dotación propia para NOC/SOC y mesa de servicios, diferenciándola de los cuatro TI del cliente. \textbf{Estado y cierre:} Dimensionamiento presente; los cuatro TI del cliente no operan solos toda la plataforma. & SD4, §4.2.6; Anexo 4-W.7. \\
\LFXfiletefila
% ID: SD4-43
\mbox{SD4-43} & \textbf{EF-22} — Sin red, despliegue ni zonas de seguridad, p. 13 & Se desarrollan cinco ambientes, topología AWS, tres vistas de red y direccionamiento con zonas de segmentación. \textbf{Estado y cierre:} Red y segmentación incorporadas documentalmente; pendientes legibilidad y cotejo de rutas, zonas y recursos de recuperación. & SD4, §§4.2.4.1–4.2.4.2; figuras 24–26; Tabla 9. \\
\LFXfiletefila
% ID: SD4-44
\mbox{SD4-44} & \textbf{EF-23} — TPS sin sustento y regímenes mezclados, p. 13 & Los cálculos separan operación normal, ventanas horarias y carga de septiembre, con su derivación en la memoria. \textbf{Estado y cierre:} Cálculos documentados; prueba no ejecutada. & SD4, §§4.2.6.2–4.2.6.3; Anexo 4-W. \\
\LFXfiletefila
% ID: SD4-45
\mbox{SD4-45} & \textbf{EF-24} — Dieciséis dimensiones ausentes o sin derivación, pp. 13–14 & Las dieciséis dimensiones se sintetizan en el anexo y se justifican mediante derivaciones y cuadros numéricos del cuerpo. \textbf{Estado y cierre:} Método documentado; cotejar suficiencia de cada dimensión y validar cotas por medición. & SD4, §§4.2.6.1–4.2.6.13; Anexo 4-W, Tabla A.30 y memoria de cálculo. \\
\LFXfiletefila
% ID: SD4-46
\mbox{SD4-46} & \textbf{EF-25} — ADR con precios, p. 13 & El registro único de ADR se presenta mediante decisiones cualitativas y técnicas. La ausencia de precios debe mantenerse en la revisión final de la oferta. \textbf{Estado y cierre:} Mantener control de ausencia de precios en toda la oferta técnica. & Anexo 4-O; control final del conjunto de la oferta técnica. \\
\LFXfiletefila
% ID: SD4-47
\mbox{SD4-47} & \textbf{EF-26} — Remisión circular sobre los 14 recintos: la página 19 remite a Seguridad §12 y esta se cita a sí misma sin enumerarlos, p. 12 & La descripción actual organiza el recinto de Talca en 4.3.1.4 y remite directamente a su distribución dibujada, sin conservar aquella remisión a Seguridad §12. \textbf{Estado y cierre:} Respuesta textual documentada; cotejar distribución y referencias en el PDF final. & SD4, §4.3.1.4, Figura 29. \\
\LFXfiletefila
% ID: SD4-48
\mbox{SD4-48} & \textbf{EF-27} — Valoración positiva: TPS derivados con método e identificación del primer cuello de botella, p. 13 & Se conserva el principio valorado: tasas con entradas, ventanas y derivaciones reproducibles, junto con el análisis del primer cuello de botella. Los valores y el candidato cambian con el diseño actual; no se presentan los números antiguos como resultados vigentes. \textbf{Estado y cierre:} Método documentado; los resultados medidos corresponden a la aceptación prevista. & SD4, §§4.1.19, 4.2.6.2–4.2.6.3 y 4.2.6.11; Anexo 4-W. \\
\LFXfiletefila
% ID: SD4-49
\mbox{SD4-49} & \textbf{EF-28} — Valoración positiva: quince ADR en formato MADR, con fecha, estado y alternativas evaluadas, p. 13 & Se conserva el registro de decisiones y sus alternativas, ampliado a 22 ADR en el diseño actual. El resumen y los estados se concentran en el anexo lógico; la retirada de precios se sigue por EF-25. \textbf{Estado y cierre:} Registro documentado; verificar vigencia y selección en cada ficha antes del cierre. & Anexo 4-O, Tablas A.18–A.19 y fichas ADR-01 a ADR-22. \\
\LFXfiletefila
% ID: SD4-50
\mbox{SD4-50} & \textbf{EL-22} — Hoja «Página en blanco a propósito» al inicio del archivo, apartado General & El cuerpo integrado comienza en la portada y no contiene esa hoja en ninguna de sus páginas; los anexos y el T-11 tampoco. \textbf{Estado y cierre:} Corregido. & SD4, anexos y T-11. \\
\LFXfiletefila
% ID: SD4-51
\mbox{SD4-51} & \textbf{EF-29} — La palabra LAFROX no aparece y se lee «Empresa proponente: Pendiente de definir», apartado General & LafroX figura en la portada y en todas las páginas del cuerpo, de los anexos y del T-11; la expresión «Pendiente de definir» no aparece. \textbf{Estado y cierre:} Corregido. & SD4, anexos y T-11, portada, encabezado y pie. \\
\LFXfiletefila
% ID: SD4-52
\mbox{SD4-52} & \textbf{EF-30} — Más de la mitad del documento son tablas de listado, contra el Formulario T-21, apartado General & El detalle se trasladó: el T-11 lista equipos y servicios, y los 23 anexos 4-A a 4-W reúnen catálogos, matrices y la memoria de cálculo. El cuerpo conserva 26 tablas técnicas breves y la declaración de IA. \textbf{Estado y cierre:} Documentado. & SD4, cuerpo; anexos 4-A a 4-W; Formulario T-11. \\
\end{longtable}
}
\Needspace{18\baselineskip}
\subsection{Subdocumento 5}
{\linespread{1}\LFXcuerpotabla
\setlength{\parskip}{0pt}
\renewcommand{\arraystretch}{1.10}
\setlength{\tabcolsep}{4pt}
\begin{longtable}{F{0.10} F{0.27} F{0.43} F{0.20}}
\caption{Observaciones y respuestas --- subdocumento 5}\\[2pt]
\LFXtopruleoriginal \cab{N.º} & \cab{Observación del Informe 1} & \cab{Respuesta} & \cab{Sección modificada} \\ \midrule
\endfirsthead
\multicolumn{4}{@{}l@{}}{\fontsize{9pt}{11pt}\selectfont\color{lafroxGris}Tabla \thetable\ --- Subdocumento 5 (continuación)}\\[3pt]
\LFXtopruleoriginal \cab{N.º} & \cab{Observación del Informe 1} & \cab{Respuesta} & \cab{Sección modificada} \\ \midrule
\endhead
\midrule
\multicolumn{4}{@{}r@{}}{\fontsize{9pt}{11pt}\selectfont\color{lafroxGris}Continúa en el folio siguiente.}\\
\endfoot
\bottomrule
\endlastfoot
% ID: SD5-01
\mbox{SD5-01} & No hay modelo de datos; 5.1 solo enumera quince bases lógicas. & Se incorporaron MR general, vistas por dominio y modelo dimensional: diez figuras en el cuerpo. & 5.1; Figuras 5.1–5.10; 5-M \\
\LFXfiletefila
% ID: SD5-02
\mbox{SD5-02} & No hay diccionario real; se posterga hasta antes de producción. & Diccionario de 121 entidades y 987 atributos, con tipos, dominios, claves, condiciones y controles. & 5-A, A.2–A.10; 5-B–5-C \\
\LFXfiletefila
% ID: SD5-03
\mbox{SD5-03} & No pueden evaluarse las consultas del proveedor y la autoridad sanitaria. & Cadena lote/unidad/evento/entrega/documento y consulta inversa; universo potencial offline y presupuesto de retiro de 85 min. & 5.1.7; Figura 5.5; 5-J \\
\LFXfiletefila
% ID: SD5-04
\mbox{SD5-04} & Motor y CAP se presentan en tabla sin analizar compromisos. & Justificación en prosa de motores, autoridad única por agregado y consistencia local con consolidación pendiente. La tabla de cinco columnas con palabras partidas se reemplazó por el desarrollo en prosa de 5.2.1; la correspondencia por dominio pasa al Anexo 5-B. & 5.2.1–5.2.3; 5-B \\
\LFXfiletefila
% ID: SD5-05
\mbox{SD5-05} & Se cita la decisión 16.1 n.º 18, inexistente. & Referencia retirada; las fuentes identifican decisiones existentes del numeral 16.1. & 5.1; fuentes de Figuras 5.1, 5.2, 5.4 y 5.6 \\
\LFXfiletefila
% ID: SD5-06
\mbox{SD5-06} & Alcance por migrar con los valores del Cap. 15 y casos de saneamiento propios del caso: lotes en texto libre, duplicados de cliente y saldos de envases. & Históricos delimitados; originales preservados, equivalencias aprobadas y defectos aislados sin inventar historia. & 5.3.1–5.3.5; 5-C; 5-L, A.25 \\
\LFXfiletefila
% ID: SD5-07
\mbox{SD5-07} & Migración sin volumen, herramientas, tiempos, ventana ni responsables. & Fuente 16,05478 GB; destino estimado 32,10956 GB; herramientas, responsables, dos ensayos, olas, corte y reversa. & 5.3; Tabla 5.2; 5-J; 5-L \\
\LFXfiletefila
% ID: SD5-08
\mbox{SD5-08} & OK: índices para retiro, partición mensual y orden en mensajería. & Se conservan; claves globales sin partición y detalle temporal con identidad completa, outbox y deduplicación. & 5.4.2–5.4.5; 5-F–5-G \\
\LFXfiletefila
% ID: SD5-09
\mbox{SD5-09} & La expresión de TPS base y crecimiento es ininteligible. & Se separan régimen 12,34 TPS, pico 14,66 y prueba 21,99; año 3 y crecimiento 3× son escenarios distintos. & 5.4.1; Tabla 5.3; 5.4.6 \\
\LFXfiletefila
% ID: SD5-10
\mbox{SD5-10} & Identidad de bodega con TTL de 8 h frente a autonomía exigida de 24 h. & Se distinguen autonomía 24 h, identidad 24 h, turnos 8/14 h y manifiesto 26 h con relevos preinscritos. & 5.2.3; Tabla 5.1; 5-G \\
\LFXfiletefila
% ID: SD5-11
\mbox{SD5-11} & OK: separación OLTP/analítica, latencias, retenciones, inmutabilidad e ISO 25012. & Se mantienen y precisan calidad, archivo, purga, retención legal y conservación de reglas interpretativas. & 5.2.5–5.2.7; 5-D–5-E; 5-H–5-I \\
\LFXfiletefila
% ID: SD5-12
\mbox{SD5-12} & OK: detecta desajuste RT-05.10/RT-16.10 y lo eleva a consulta. & Correspondencia y consulta documentadas; plazos íntegros conservados. No se presume envío acreditado ni respuesta oficial. & 5.2.7; 5-D, A.14 \\
\LFXfiletefila
% ID: SD5-13
\mbox{SD5-13} & Faltan esquema estrella dibujado, tableros, destinatarios e indicadores definidos. & Modelo dimensional con granos y claves; fórmulas, audiencia, latencia y linaje. El costo de servir identifica combustible, peajes, conducción, descarga e indirectos, con prorrateos versionados y conciliación con el ERP. & 5.1.11; Figura 5.7; 5.2.6; 5-H, Tablas A.18–A.19 \\
\LFXfiletefila
% ID: SD5-14
\mbox{SD5-14} & RT-16.08 se atribuye a tres canales; corresponde RT-16.20. & RT-16.20 identifica canales y RT-16.08 consulta/exportación de auditoría. & 5.1.9; 5-B–5-C; A.8 \\
\LFXfiletefila
% ID: SD5-15
\mbox{SD5-15} & Timbre diferido y guía ausente contradicen ERP único emisor y acompañamiento del traslado. & ERP emite antes del traslado; cambios requieren nueva guía válida. POD no sustituye documento fiscal. & 5.1.8; 5.3.6; 5-J, AL-DTE-01 \\
\LFXfiletefila
% ID: SD5-16
\mbox{SD5-16} & OK: evidencia por receptor; falta mecanismo de acuse legal. & POD, acuse SII y recibo se diferencian; XML ERP/EDI o constancia impresa según perfil, conforme al Anexo 4-U. & 5.1.8; A.7; 5-C; 5-J \\
\LFXfiletefila
% ID: SD5-17
\mbox{SD5-17} & Resaltado negro ilegible en páginas 3, 10, 13 y 14. & Retirados los bloques de código con fondo negro; contenido integrado en la plantilla de texto, tablas y figuras. & Cuerpo y anexos \\
\LFXfiletefila
% ID: SD5-18
\mbox{SD5-18} & Metadatos finales citan documentos inexistentes. & Línea de metadatos retirada; referencias identifican Bases, propuesta y fuentes técnicas concretas, con T-12 adscrito al SD3, T-11 al SD4, T-14 incluido en el anexo y versión EPCIS alineada. & Referencias del cuerpo y anexo \\
\LFXfiletefila
% ID: SD5-19
\mbox{SD5-19} & Doble numeración 5.1/S5.1 y códigos internos sin explicación. & Numeración 5.1–5.4 y anexos 5-A–5-M; glosario y leyendas definen los identificadores necesarios. & Introducción; 5-A, glosario y A.1 \\
\LFXfiletefila
% ID: SD5-20
\mbox{SD5-20} & Siglas sin definir, incluida la frase sobre FID. & FID y la frase observada se retiraron; glosario común y desarrollo de CAP en el texto. & Introducción; 5.2.1; 5-A \\
\LFXfiletefila
% ID: SD5-21
\mbox{SD5-21} & General: ningún archivo tiene identidad de empresa; las portadas incluyen datos ajenos a la licitación y se rotulan «Formulario T-22». & La portada identifica a LafroX, al mandante, al proponente y a su representante, y se rotula «Subdocumento 5 · Formulario T-7», sin datos del curso. & Portada \\
\LFXfiletefila
% ID: SD5-22
\mbox{SD5-22} & General: el archivo abre con una hoja que sólo dice «Página en blanco a propósito». & La hoja no aparece en el cuerpo ni en los anexos. & Inicio del documento \\
\LFXfiletefila
% ID: SD5-23
\mbox{SD5-23} & General: Arial y tamaño A4, cuerpo que aparenta ser inferior a 11 puntos, y ningún índice, folio ni firma. & El cuerpo tiene índice general, de tablas y de figuras, y los anexos su índice. El tamaño de página, la tipografía, el folio y la media firma se cotejan en el PDF compilado. & Índices; formato de página, pie y firma por cotejar en el PDF \\
\LFXfiletefila
% ID: SD5-24
\mbox{SD5-24} & Transversal: indicios de IA (IFEP01.26-A1); desde el Informe 2, cualquiera deja el subdocumento con puntaje 0. & No quedan marcadores, notas de proceso ni metadatos del asistente (filas 17 a 20). La declaración de uso de IA identifica a Bastian Trejo como revisor de cada sección y anexo. & Declaración de uso de IA, Tabla 5.4 \\
\end{longtable}
}
\Needspace{18\baselineskip}
\subsection{Subdocumento 13}
{\linespread{1}\LFXcuerpotabla
\setlength{\parskip}{0pt}
\renewcommand{\arraystretch}{1.10}
\setlength{\tabcolsep}{4pt}
\begin{longtable}{F{0.10} F{0.27} F{0.43} F{0.20}}
\caption{Observaciones y respuestas --- subdocumento 13}\\[2pt]
\LFXtopruleoriginal \cab{N.º} & \cab{Observación del Informe 1} & \cab{Respuesta} & \cab{Sección modificada} \\ \midrule
\endfirsthead
\multicolumn{4}{@{}l@{}}{\fontsize{9pt}{11pt}\selectfont\color{lafroxGris}Tabla \thetable\ --- Subdocumento 13 (continuación)}\\[3pt]
\LFXtopruleoriginal \cab{N.º} & \cab{Observación del Informe 1} & \cab{Respuesta} & \cab{Sección modificada} \\ \midrule
\endhead
\midrule
\multicolumn{4}{@{}r@{}}{\fontsize{9pt}{11pt}\selectfont\color{lafroxGris}Continúa en el folio siguiente.}\\
\endfoot
\bottomrule
\endlastfoot
% ID: SD13-01
\mbox{SD13-01} & Innovación 1: «es el criterio de aceptación N° 1 del caso y el estándar exigido por el RT-05.23». & Se reemplaza por el seguimiento de vencimiento posentrega en el local (INN-01). La trazabilidad de lote hasta el cliente queda como alcance obligatorio (RF-02.10 y RF-01.03) y la innovación actúa en el local, después de la entrega. & 13.1; 13.1.2 (Tabla 13.2); Formulario T-19, ficha 1 \\
\LFXfiletefila
% ID: SD13-02
\mbox{SD13-02} & Innovación 3: «es el RT-03.10 a RT-03.13 y las restricciones N° 2 y N° 3». & Se reemplaza por la vida útil remanente por historia térmica del lote (INN-03). La operación sin conexión y el registro continuo de temperatura quedan como alcance obligatorio; la innovación agrega la estimación por lote. & 13.3; 13.3.2 (Tabla 13.12); Formulario T-19, ficha 3 \\
\LFXfiletefila
% ID: SD13-03
\mbox{SD13-03} & Innovación 2: el picking por olas y el FEFO son estándar; «lo propio sería la secuenciación térmica, que no se diseña». & Se reemplaza por la reproducción de incidentes de terreno (INN-02), con su circuito de cuatro pasos, su contrato de evidencia y la explicación de por qué no es práctica estándar. El FEFO queda como alcance obligatorio (RF-05.05). & 13.2; 13.2.2 (Tabla 13.7); Anexo 13.A \\
\LFXfiletefila
% ID: SD13-04
\mbox{SD13-04} & Innovación 4: «no coincide con la restricción N° 9 y no se refleja en la estructura de costos ni en el flujo de caja». & Se reemplaza por el tramo variable de la Operación ligado al costo de servir (INN-04). El acompañamiento se ofrece como servicio costeado, según la restricción 9, y el flujo de caja separa costo, ingreso fijo, ingreso contingente y beneficio del CLIENTE por mes. & 13.4.1; 13.4.5 (Tabla 13.20); Formulario T-19, ficha 4 \\
\LFXfiletefila
% ID: SD13-05
\mbox{SD13-05} & Innovación 5: supone WhatsApp en un cliente sin internet (restricción N° 5) y no existe en el catálogo del Subdocumento 3 ni en los módulos del 4.1 (RT-26.01). & Se reemplaza por la hoja de negocio del almacenero (INN-05), entregada en la visita en la pantalla del preventista o impresa en la entrega siguiente, sin dispositivo ni conexión del cliente. Usa M3, M6 y M10, la aplicación C-01 y la impresora de cabina ya ofertadas. & 13.5.2; 13.5.4; Anexo 13.B \\
\LFXfiletefila
% ID: SD13-06
\mbox{SD13-06} & Falta el nivel de madurez con la escala utilizada (RT-26.03). & Cada innovación declara su nivel en la escala TRL de la Comisión Europea (2016), con el nivel de partida, los niveles objetivo con su mes y las fuentes que los sustentan. & 13.1.3 a 13.5.3; Formulario T-19, campo 5 \\
\LFXfiletefila
% ID: SD13-07
\mbox{SD13-07} & Falta el diseño de incorporación: capa y componentes de la arquitectura, paquetes de la EDT y mes del cronograma (RT-26.01 y RT-26.02). & Cada innovación declara capas, componentes, módulos, contrato de datos, paquetes del Formulario T-14 y mes de materialización. & 13.1.4 a 13.5.4; Anexos 13.A y 13.B; Formulario T-19, campos 7 a 9 \\
\LFXfiletefila
% ID: SD13-08
\mbox{SD13-08} & Falta el impacto económico: inversión, costo operacional y beneficio (RT-26.05). & Cada innovación ubica sus rubros por clase y mes en el flujo de caja, sin montos (Art. 50.2). & 13.1.5 a 13.5.5; Anexo 13.B; Formulario T-19, campos 10 a 12 \\
\LFXfiletefila
% ID: SD13-09
\mbox{SD13-09} & Falta el riesgo de adopción con probabilidad, impacto, mitigación y contingencia (RT-26.04). & Cada innovación remite a su ficha R8-25 a R8-29 del Anexo 8.A, con probabilidad, impacto y exposición en la escala del Capítulo 8, mitigación y contingencia. & 13.1.7 a 13.5.7; Anexo 13.C (Tabla 13.C.2); Formulario T-19, campos 15 a 17 \\
\LFXfiletefila
% ID: SD13-10
\mbox{SD13-10} & «“Investigación adicional: Sí” en las cinco equivale a declarar que la innovación no está diseñada». & Cada innovación declara qué falta investigar, en qué paquete y meses se resuelve y con qué indicador se valida. & 13.1.3 a 13.5.3; Anexo 13.B (Tabla 13.B.2) \\
\LFXfiletefila
% ID: SD13-11
\mbox{SD13-11} & Los indicadores no tienen momento de medición ni línea base; la innovación 5 compromete un valor al mes 24 de Operación, sin medición intermedia. & Cada indicador tiene línea base, meta y momento de medición. Donde no existe historia se declara «Sin medición» y no se inventa un valor. La innovación 5 mide en los meses 13 y 14, 27 y 33. & Tablas 13.5, 13.10, 13.15, 13.21 y 13.26; Anexo 13.C (Tabla 13.C.1) \\
\LFXfiletefila
% ID: SD13-12
\mbox{SD13-12} & Crítico: el Formulario T-19 no fue entregado. & Se presenta el Formulario T-19 en archivo propio, con una ficha por innovación y sus 17 campos. & Formulario T-19 \\
\LFXfiletefila
% ID: SD13-13
\mbox{SD13-13} & Crítico: la innovación 5 trae precios en dólares (Art. 50.2). & Ningún monto en el capítulo ni en el formulario. El impacto se expresa por rubro, clase y mes, y la valorización va en la Oferta Económica, Entregable 2, ítem 2.7. & 13.1.5 a 13.5.5; Formulario T-19, campos 10 a 12 \\
\LFXfiletefila
% ID: SD13-14
\mbox{SD13-14} & Las fuentes no se citan dentro del texto; la versión 1.2 de EPCIS está superada por la 2.0 (2022). & Las fuentes se citan en APA 7.ª en el punto de uso, con correspondencia uno a uno con la lista. EPCIS 2.0 (GS1 AISBL, 2022) sustenta la tecnología y la madurez de la innovación 3. & 13.1.3 a 13.5.3; 13.3.2; 13.3.3; Referencias \\
\LFXfiletefila
% ID: SD13-15
\mbox{SD13-15} & «Siete páginas sin una figura: ninguna innovación se ubica en un diagrama de la arquitectura». & La Figura 13.1 ubica las cinco innovaciones sobre las capas de la arquitectura, y las Figuras 13.2 a 13.6 recorren cada una. Cada figura se cita antes y se explica después. & Introducción (Figura 13.1); 13.1.4 a 13.5.4 \\
\LFXfiletefila
% ID: SD13-16
\mbox{SD13-16} & OK: el anexo descarta bien lo que es requisito y no innovación: ruteo automático, costo de servir y registro continuo de temperatura. & Se mantiene el criterio. Cada innovación separa en una tabla lo que ya es alcance obligatorio de lo que agrega; el registro continuo de temperatura (criterio 3) y el costo de servir (criterio 11 y RF-11.08) se declaran alcance obligatorio y no innovación. & Tablas 13.2, 13.7, 13.12, 13.17 y 13.23; 13.3; 13.4.1 \\
\LFXfiletefila
% ID: SD13-17
\mbox{SD13-17} & General: ningún archivo tiene logotipo, plantilla, paleta, encabezado ni pie de página. & Plantilla única de LafroX, con isotipo, encabezado y pie en todas las páginas. & Subdocumento, anexos y Formulario T-19 \\
\LFXfiletefila
% ID: SD13-18
\mbox{SD13-18} & General: las portadas incluyen datos ajenos a la licitación y están rotuladas «Formulario T-22», que es un formulario del CLIENTE. & Portada con la identidad de LafroX, sin datos ajenos a la licitación ni rótulo de formularios del CLIENTE. & Portada \\
\LFXfiletefila
% ID: SD13-19
\mbox{SD13-19} & General: los Subdocumentos 1, 2, 3, 5 y 13 en Arial y tamaño A4; cuerpo que aparenta ser inferior a 11 puntos. & Tamaño carta y cuerpo de 11 puntos (Art. 40.4). & Subdocumento, anexos y Formulario T-19 \\
\LFXfiletefila
% ID: SD13-20
\mbox{SD13-20} & General: el archivo abre con una hoja que sólo dice «Página en blanco a propósito», con errata. & Se elimina la hoja. & Inicio del documento \\
\LFXfiletefila
% ID: SD13-21
\mbox{SD13-21} & General: ningún archivo tiene índice, numeración de páginas ni firma (Art. 40° y 53°). & Índice detallado hasta el tercer nivel, folio correlativo y zona de media firma en el extremo inferior derecho de cada página. & Índice; pie de página \\
\LFXfiletefila
% ID: SD13-22
\mbox{SD13-22} & Transversal: marcadores, instrucciones y notas de proceso dentro del entregable; nombres de archivo como fuente. & El capítulo no contiene marcadores, notas de proceso ni nombres de archivo como fuente, y cierra con la Declaración de uso de IA por sección, anexo y formulario (Aclaración §7.2). & Subdocumento completo; Declaración de uso de IA \\
\LFXfiletefila
% ID: SD13-23
\mbox{SD13-23} & Transversal: la propuesta se contradice entre subdocumentos en decisiones de diseño. & Las innovaciones usan los componentes, módulos, meses y responsables del Subdocumento 1, del Capítulo 4, del Formulario T-14, del Anexo 7.D y del Anexo 8.A. & 13.1.4 a 13.5.4; Anexo 13.B \\
\LFXfiletefila
% ID: SD13-24
\mbox{SD13-24} & Transversal: no es posible seguir un requerimiento desde el problema hasta un componente de la arquitectura. & Cada innovación separa lo que ya es alcance obligatorio con sus RF y RT, y el Anexo 13.B la traza a capas, componentes, módulos, paquetes, riesgos y flujo de caja. & Tablas 13.2, 13.7, 13.12, 13.17 y 13.23; Anexo 13.B \\
\end{longtable}
}
\clearpage
\subsection{Observación transversal}
{\linespread{1}\LFXcuerpotabla
\setlength{\parskip}{0pt}
\renewcommand{\arraystretch}{1.10}
\setlength{\tabcolsep}{4pt}
\begin{longtable}{F{0.10} F{0.27} F{0.43} F{0.20}}
\caption{Observaciones y respuestas --- observación transversal}\\[2pt]
\LFXtopruleoriginal \cab{N.º} & \cab{Observación del Informe 1} & \cab{Respuesta} & \cab{Sección modificada} \\ \midrule
\endfirsthead
\multicolumn{4}{@{}l@{}}{\fontsize{9pt}{11pt}\selectfont\color{lafroxGris}Tabla \thetable\ --- Observación transversal (continuación)}\\[3pt]
\LFXtopruleoriginal \cab{N.º} & \cab{Observación del Informe 1} & \cab{Respuesta} & \cab{Sección modificada} \\ \midrule
\endhead
\midrule
\multicolumn{4}{@{}r@{}}{\fontsize{9pt}{11pt}\selectfont\color{lafroxGris}Continúa en el folio siguiente.}\\
\endfoot
\bottomrule
\endlastfoot
% ID: TR-01
\mbox{TR-01} & Transversal: fundamentación ingenieril. Hay decisiones correctas y propias del caso, pero enterradas en tablas: sin diagramas explicados, sin cálculos a la vista y sin análisis en prosa. & \textbf{Documentado.} Las decisiones se presentan con figuras explicadas en el texto: Figura 1.1 en el SD1, Figuras 2.1–2.6 en el SD2, Figuras 3.1–3.8 en el SD3, figuras 1–29 en el SD4, Figuras 5.1–5.10 en el SD5 y Figuras 13.1–13.6 en el SD13. Los cálculos quedan a la vista: dotación del NOC en el SD1 (4×7+4=32), reentregas en el SD2 (1.400 × 0,042 = 58,8 ≈ 59 al día), memoria de cálculo del SD4 (§4.2.6 y Anexo 4-W) y volumen de migración y escenarios de carga del SD5 (Tablas 5.2 y 5.3). Las decisiones se analizan en prosa: reglas sensibles del SD3 (§3.2.4), comparación de alternativas y 22 ADR del SD4 (§4.1.13 y Anexo 4-O) y compromisos de persistencia del SD5 (§5.2). & SD1–SD5 y SD13; filas SD1-04, SD1-16, SD2-13, SD2-31, SD3-07, SD3-35, SD4-09, SD4-45, SD5-01, SD5-07, SD5-09 y SD13-15. \\
\end{longtable}
}
\Needspace{6\baselineskip}
\section{Notas por subdocumento}
\Needspace{6\baselineskip}
\subsection{Subdocumento 1}
La tabla relaciona las observaciones del Informe 1 con sus respuestas y las secciones correspondientes del Subdocumento 1, sus anexos y formularios asociados. «Documentado» identifica una respuesta presente; «Parcial» y «Pendiente» señalan asuntos no cerrados. Estos estados no equivalen a aceptación del CLIENTE.

\Needspace{6\baselineskip}
\subsection{Subdocumento 2}
La tabla relaciona las observaciones del Informe 1 con sus respuestas y las secciones correspondientes del Subdocumento 2, sus anexos y formularios asociados. «Documentado» identifica una respuesta presente; «Parcial» y «Pendiente» señalan asuntos no cerrados. Estos estados no equivalen a aceptación del CLIENTE.

\Needspace{6\baselineskip}
\subsection{Subdocumento 3}
La tabla relaciona las observaciones del Informe 1 con sus respuestas y las secciones correspondientes del Subdocumento 3, sus anexos y formularios asociados. «Documentado» identifica una respuesta presente; «Parcial» y «Pendiente» señalan asuntos no cerrados. Estos estados no equivalen a aceptación del CLIENTE.

\Needspace{6\baselineskip}
\subsection{Subdocumento 4}
\Needspace{6\baselineskip}
\subsubsection{Fuente y alcance}
Esta tabla responde a los apartados \textbf{SUBDOCUMENTO 4.1} y \textbf{SUBDOCUMENTO 4.2} de G02-C02-LOGISTICA-LAFROX-Revision Informe 1.pdf, páginas 9–14. Se excluyen el contenido del SD3 anterior al inicio del apartado 4.1 en la página 9 y el SD5 posterior al cierre del apartado 4.2 en la página 14. Las instrucciones del PDF se tratan como observaciones evaluativas.

Las observaciones se sintetizan, sin pretender ser citas literales ni reproducir una numeración que el PDF no contiene. Se conservan EL-01–EL-21 y EF-01–EF-25 como identificadores del registro previo, y se añaden EF-26–EF-28 para tres puntos explícitos del PDF que no tenían respuesta propia. Se añaden también EL-22, EF-29 y EF-30 para los puntos del apartado General del PDF que alcanzan al Subdocumento 4; su primera columna indica ese apartado en lugar de una página. Algunas filas agrupan observaciones del mismo tema; los puntos positivos también se conservan. Las exigencias de «Qué se espera en el Informe 2» se vinculan a los temas correspondientes, incluida la presentación con índice, folio y firma en EF-01.

Las páginas indicadas en las observaciones corresponden al \textbf{PDF de revisión}, no al entregable actual. La columna «Sección modificada» localiza la respuesta en SD4, anexos y T-11. Las referencias a SD3 o T-12 se mantienen únicamente cuando la observación del SD4 requiere comprobar consistencia o trazabilidad. No se incluye una auditoría independiente de esos documentos.

«Documentado» significa que existe una respuesta de diseño o contenido identificada. No acredita aprobación del evaluador, exactitud de cada afirmación, pruebas ejecutadas ni conformidad final. La cobertura de observaciones no equivale a una acreditación integral de los entregables.

La tabla relaciona cada observación con su respuesta, condición de cierre y ubicación documental.

\Needspace{6\baselineskip}
\subsubsection{Qué falta para acreditar el cierre}
La cobertura del PDF adjunto queda registrada en la tabla, incluidos los tres puntos añadidos. No quedan identificadas observaciones del SD4 excluidas deliberadamente. Sin embargo, la cobertura de una tabla de respuestas no demuestra que todos los reparos del entregable estén cerrados. Siguen requiriendo comprobación los siguientes puntos, vinculados a las observaciones originales:

\begin{itemize}
\item \begin{samepage}
\textbf{Trazabilidad individual y consistencia entre capítulos:} demostrar la suficiencia de la cadena requisito → módulo → capa → componente físico y cotejar etapas con SD3, no solo la existencia de matrices (EL-04, EL-17 y EF-03/24).
\end{samepage}
\item \begin{samepage}
\textbf{Figuras y presentación final:} cotejar actores, etiquetas, recorridos y emplazamientos en los dibujos; comprobar legibilidad impresa, índice, folios y firma de los archivos finales (EL-01/03/20/21 y EF-01/04/05/13/22/26).
\end{samepage}
\item \begin{samepage}
\textbf{Identidad y operación sin enlace:} verificar coherencia del permiso local, relevos, gateway, bloqueo térmico, funciones no disponibles y guía emitida por ERP antes de la salida (EL-14/15/16 y EF-18/20). Los ensayos previstos pertenecen a la aceptación del proyecto; no se exige inventar resultados para responder al Informe 1.
\end{samepage}
\item \begin{samepage}
\textbf{Recuperación y residencia:} cerrar la condición de aprobación de la región secundaria y demostrar que la alternativa comprometida conserva el RTO requerido; la sola declaración del objetivo no resuelve la observación (EF-15).
\end{samepage}
\item \begin{samepage}
\textbf{Implementos y capacidad:} cotejar cantidades por tipología, resistencia a caída, ciclo de vida, congelado, impresoras/POS, parque de 28 camiones, estudio inalámbrico, disponibilidad, almacenamiento y las dieciséis estimaciones con sus entradas y derivaciones (EF-08–EF-14, EF-19 y EF-23/24). Las mediciones futuras se distinguen de los cálculos de oferta.
\end{samepage}
\item \begin{samepage}
\textbf{Limpieza de la oferta:} comprobar ausencia de comentarios al redactor, referencias históricas sin destino, precios y desarrollo de BI ajeno al propósito arquitectónico (EL-20 y EF-02/06/25).
\end{samepage}
\end{itemize}
Estos son límites de cierre documental y de evidencia; no se afirma que todos sean defectos aún presentes ni que la ejecución de pruebas futuras deba haber ocurrido antes de presentar la propuesta.

\Needspace{6\baselineskip}
\subsection{Subdocumento 5}
Respuestas documentales del cuerpo y los anexos vigentes, conforme al Art. 46. La tabla contiene las veinte observaciones del apartado del Subdocumento 5 en el Informe 1 y cuatro observaciones generales que lo alcanzaban, cada una con su respuesta y sección modificada. Los ensayos siguen programados para el proyecto; no se presume respuesta oficial a la consulta de códigos.

\Needspace{6\baselineskip}
\subsection{Subdocumento 13}
El Art. 46 de las Bases Administrativas pide resolver en forma explícita las observaciones de la instancia anterior, con una tabla de trazabilidad observación–respuesta–sección modificada (Distribuidora Puelche S.A., 2026a, Art. 46, p. 28). Las filas SD13-01 a SD13-16 responden las dieciséis observaciones al Capítulo 13, incluida la valoración positiva; las filas SD13-17 a SD13-24 responden las ocho observaciones generales y transversales aplicables.

Filas SD13-01 a SD13-16: Observaciones al Capítulo 13 de la revisión del Informe 1. Fuente: elaboración a partir de las observaciones de la instancia anterior (Distribuidora Puelche S.A., 2026a, Art. 46, p. 28).

Las cinco innovaciones del Informe 1 se reemplazaron; por eso las filas 1 a 5 responden con una innovación nueva y no con un ajuste de la anterior. En la innovación 3, el filtro del Art. 30° de las Bases Administrativas (p. 20) se acredita con el alcance descrito en 13.3.2, que distingue técnicamente la incorporación propuesta de las funcionalidades ya exigidas, además de separarla del registro continuo obligatorio.

Filas SD13-17 a SD13-24: Observaciones generales y transversales aplicables al Subdocumento 13. Fuente: elaboración a partir de las observaciones de la instancia anterior (Distribuidora Puelche S.A., 2026a, Art. 46, p. 28).

Las filas 17 a 24 recogen las observaciones de forma y de coherencia que la revisión dirigió a toda la entrega; aquí se responde sólo lo que corresponde a este subdocumento, sus anexos y su formulario.

\microtypesetup{protrusion=false}
\Needspace{6\baselineskip}
\section{Referencias}
\begin{itemize}
\item \begin{samepage}
Revisión del Informe 1, observaciones generales y apartados «SUBDOCUMENTO 1», «SUBDOCUMENTO 2», «SUBDOCUMENTO 3», «SUBDOCUMENTO 4.1», «SUBDOCUMENTO 4.2», «SUBDOCUMENTO 5» y «SUBDOCUMENTO 13».
\end{samepage}
\item \begin{samepage}
Subdocumento 1.
\end{samepage}
\item \begin{samepage}
Formulario T-6.
\end{samepage}
\item \begin{samepage}
Subdocumento 2.
\end{samepage}
\item \begin{samepage}
Anexos del Subdocumento 2.
\end{samepage}
\item \begin{samepage}
Subdocumento 3.
\end{samepage}
\item \begin{samepage}
Anexos del Subdocumento 3.
\end{samepage}
\item \begin{samepage}
Formulario T-12.
\end{samepage}
\item \begin{samepage}
Subdocumento 4.
\end{samepage}
\item \begin{samepage}
Anexos del Subdocumento 4.
\end{samepage}
\item \begin{samepage}
Formulario T-11.
\end{samepage}
\item \begin{samepage}
Subdocumento 5.
\end{samepage}
\item \begin{samepage}
Anexos del Subdocumento 5.
\end{samepage}
\item \begin{samepage}
Subdocumento 13.
\end{samepage}
\item \begin{samepage}
Anexos del Subdocumento 13.
\end{samepage}
\item \begin{samepage}
Matriz de cambios del Informe 1 al Informe 2 del Subdocumento 13.
\end{samepage}
\item \begin{samepage}
Formulario T-19.
\end{samepage}
\item \begin{samepage}
Bases Administrativas, artículo 46° y Formulario T-22.
\end{samepage}
\item \begin{samepage}
Aclaraciones de la licitación.
\end{samepage}
\end{itemize}
\Needspace{5\baselineskip}
Referencias citadas en las filas del Subdocumento 13:

Las citas a las Bases Administrativas y a las Bases Técnicas Transversales identifican el artículo o capítulo y la página impresa del PDF original, además del código RT cuando corresponde. Las referencias al caso se localizan por capítulo, sección o entrevista, sin exigir paginación. Las Aclaraciones se identifican por sección en la fuente disponible.

\begin{itemize}
\item \begin{samepage}
Distribuidora Puelche S.A. (2026a). \emph{Bases Administrativas de Licitación N.º TFEP-01/2026: Contratación de Solución Integral de Software y Servicios de Operación}, artículos 28 y 30, pp. 19–20, y artículo 46, p. 28.

\end{samepage}
\item \begin{samepage}
Distribuidora Puelche S.A. (2026b). \emph{Bases Técnicas Transversales de Licitación N.º TFEP-01/2026}, capítulo 26, RT-26.01 (p. 44), RT-26.07 y RT-26.08 (p. 45).

\end{samepage}
\item \begin{samepage}
Distribuidora Puelche S.A. (2026d). \emph{Aclaraciones de la Licitación N.º TFEP-01/2026}, secciones 7.2 y 8.

\end{samepage}
\item \begin{samepage}
European Commission. (2016). Technology readiness levels (TRL). En \emph{Horizon 2020 work programme 2016–2017: General annexes} (Anexo G). \url{https://ec.europa.eu/research/participants/data/ref/h2020/other/wp/2016_2017/annexes/h2020-wp1617-annex-g-trl_en.pdf}

\end{samepage}
\item \begin{samepage}
GS1 AISBL. (2022). \emph{EPC Information Services (EPCIS) standard} (versión 2.0). \url{https://ref.gs1.org/standards/epcis/2.0.0/}

\end{samepage}
\end{itemize}

\end{document}
